\chapter{Propagation and Evolution}
\label{chap:propagation}
\markboth{Propagation and Evolution}{Propagation and Evolution}

\begin{remark}[Positioning of this chapter and strict premises]\label{rem:propagation_chapter_positioning}
This chapter takes \textbf{the propagation of light} as its basic prototype and establishes a unified mathematical framework for propagation and evolution. All definitions, propositions, lemmas, theorems and corollaries in the chapter are bounded by the following two premises:

\begin{enumerate}[label=(\alph*)]
    \item \textbf{Linearity premise}: the field quantities satisfy the \textbf{linear} wave equation (the degenerate form, in a source-free region, of the vacuum Maxwell equations, \cref{thm:electromagnetic_wave_equation} of the `Periodic Structure of Electromagnetic Systems' chapter), or more generally satisfy a \textbf{linear} partial differential equation with constant coefficients (or with slowly varying coefficients);
    \item \textbf{Slowly varying medium premise}: the refractive index $n(\mathbf{r})$ of the medium is a $C^2$ smooth function, and its relative change over one wavelength scale is far less than $1$: $|\nabla n|\lambda \ll n$.
\end{enumerate}

Any statement going beyond the above premises---nonlinear media, strong-field effects, quantum optics, divergence of the field strength on a caustic surface, localisation in random media---is invariably marked as a \textbf{note} or an \textbf{open problem}, and \textbf{is not used as a proved conclusion}.

Positioning of the seven themes of this chapter: translation---spatial propagation is generated by the action of the translation group $\mathbb{R}^n$ (\cref{def:propagation_operator}), whilst time evolution is generated by the evolution semigroup (\cref{def:evolution_semigroup}); scaling---the refractive index of the medium gives rise to the wavelength scaling $\lambda = \lambda_0/n$ (\cref{prop:medium_wavelength_scaling}); rotation---the complex phase representation and the phase operator (\cref{prop:phase_propagation}); conservation---light propagation in vacuum preserves energy (\cref{thm:vacuum_energy_conservation}); symmetry---vacuum light propagation is invariant under spatial translation (\cref{prop:propagation_commutes_translation}); closure---Bloch waves in a periodic medium (\cref{thm:bloch_floquet}); period---the frequency-domain support of the dispersion relation (\cref{prop:dispersion_period}).
\end{remark}

\begin{figure}[htbp]
\centering
\begin{tikzpicture}[
  node distance=8mm and 10mm,
  box/.style={draw, rounded corners, align=center, minimum width=2.4cm, minimum height=0.8cm, fill=blue!5},
  core/.style={draw, very thick, rounded corners, align=center, minimum width=3.2cm, minimum height=1cm, fill=orange!15},
  arrow/.style={-{Stealth[length=2mm]}, thick}
]
\node[core] (core) {Propagation and evolution\\Unified mathematical framework};
\node[box, above left=of core] (trans) {Translation\\Propagation operator / evolution semigroup};
\node[box, above=of core] (scale) {Scaling\\Wavelength scaling / learning rate};
\node[box, above right=of core] (rot) {Rotation\\Complex phase / CEAF};
\node[box, left=of core] (cons) {Conservation\\Energy / flux / natural parameters};
\node[box, right=of core] (sym) {Symmetry\\Translation invariance / Bloch};
\node[box, below left=of core] (close) {Closure\\Periodic closure / band gap / convergence};
\node[box, below right=of core] (period) {Period\\Dispersion period / annealing period};
\foreach \x in {trans,scale,rot,cons,sym,close,period}
  \draw[arrow] (core) -- (\x);
\end{tikzpicture}
\caption{Positioning of the seven themes of this chapter: translation, scaling, rotation, conservation, symmetry, closure and period.}
\label{fig:chapter_map}
\end{figure}

\section{Ray trajectories: from Fermat's principle to the eikonal equation}\label{sec:ray_trajectory}

\begin{remark}[Positioning of this section]\label{rem:ray_section_positioning}
This section starts from Fermat's principle and rigorously derives the mathematical description of ray trajectories. Fermat's principle is the fundamental variational principle of geometrical optics, and making it rigorous requires the first variation of the optical-path functional. This section does not assume that the reader already knows the calculus of variations: every derivation starts from the definitions, and wherever integration by parts is involved, the treatment of the boundary terms is written out explicitly.
\end{remark}

\begin{definition}[Optical-path functional]\label{def:optical_path_functional}
Let the refractive index $n:\mathbb{R}^n\to\mathbb{R}^+$ be a $C^2$ smooth function and let $\gamma:[s_0,s_1]\to\mathbb{R}^n$ be a $C^2$ curve parametrised by arc length $s$ (that is, $|\dot\gamma(s)|=1$). The \textbf{optical-path functional} on $\gamma$ is defined by
\[
L[\gamma] := \int_{s_0}^{s_1} n(\gamma(s))\,ds.
\]
\end{definition}

\begin{definition}[Fermat's principle]\label{def:fermat_principle}
\textbf{Fermat's principle}: a light ray is a first-order stationary curve of the optical-path functional $L[\gamma]$ among all $C^2$ curves joining the two endpoints $\gamma(s_0)=\mathbf{r}_0$, $\gamma(s_1)=\mathbf{r}_1$.
\end{definition}

\begin{figure}[htbp]
\centering
\begin{tikzpicture}[>=Stealth, thick]
\coordinate (A) at (0,0);
\coordinate (B) at (5,1.2);
\draw[fill=gray!10] (0,-0.6) .. controls (2,0.4) and (3,0.8) .. (5,0.6) .. controls (5.5,0.4) and (5.5,-0.2) .. (5,-0.6) -- cycle;
\node at (2.5,-0.35) {Refractive index $n(\mathbf{r})$};
\draw[blue, very thick] (A) .. controls (1.8,1.6) and (3.2,-0.2) .. (B);
\node[blue, above] at (1.4,1.15) {True ray $\gamma(s)$};
\draw[red, dashed] (A) .. controls (2.0,2.4) and (3.0,1.8) .. (B);
\node[red, above] at (2.8,2.0) {Variation curve $\gamma_\varepsilon$};
\fill (A) circle (1.5pt) node[left] {$\mathbf{r}_0$};
\fill (B) circle (1.5pt) node[right] {$\mathbf{r}_1$};
\draw[<->] (1.2,0.2) -- (1.6,0.8);
\node (hlabel) at (-1.2,0.55) {$\mathbf{h}(s)$};
\draw[dashed, ->] (hlabel) -- (1.6,0.8);
\node at (2.5,-1.3) {$\displaystyle L[\gamma]=\int_{s_0}^{s_1} n(\gamma(s))\,ds$};
\end{tikzpicture}
\caption{Fermat's principle: variation with fixed endpoints.}
\label{fig:fermat_variation}
\end{figure}

\begin{lemma}[First variation of the optical-path functional]\label{lem:optical_path_variation}
Let $\gamma_\varepsilon(s) = \gamma(s) + \varepsilon\mathbf{h}(s)$ be a family of variation curves for $L[\gamma]$, where $\mathbf{h}\in C^2([s_0,s_1],\mathbb{R}^n)$ satisfies $\mathbf{h}(s_0)=\mathbf{h}(s_1)=\mathbf{0}$ (a variation with fixed endpoints). Then
\[
\left.\frac{d}{d\varepsilon}\right|_{\varepsilon=0} L[\gamma_\varepsilon]
= \int_{s_0}^{s_1}\Bigl[\nabla n(\gamma(s))\cdot\mathbf{h}(s)\Bigr]\,ds
+ \int_{s_0}^{s_1}n(\gamma(s))\frac{d}{ds}\Bigl[\frac{\dot\gamma(s)\cdot\dot{\mathbf{h}}(s)}{|\dot\gamma(s)|}\Bigr]_{\varepsilon=0}\,ds.
\]
Under the arc-length parametrisation $|\dot\gamma|=1$, the second term simplifies to $\int_{s_0}^{s_1}n(\gamma(s))\,\dot\gamma(s)\cdot\dot{\mathbf{h}}(s)\,ds$.
\end{lemma}

\begin{proof}
Expand $L[\gamma_\varepsilon]$ by definition and differentiate directly. Write $\mathbf{v}_\varepsilon(s) := \dot\gamma_\varepsilon(s)$; then $|\mathbf{v}_\varepsilon(s)| = |\dot\gamma(s) + \varepsilon\dot{\mathbf{h}}(s)|$. Since $|\dot\gamma|=1$,
\[
|\mathbf{v}_\varepsilon| = \sqrt{1 + 2\varepsilon\dot\gamma\cdot\dot{\mathbf{h}} + \varepsilon^2|\dot{\mathbf{h}}|^2} = 1 + \varepsilon\dot\gamma\cdot\dot{\mathbf{h}} + O(\varepsilon^2).
\]
Substituting into $L[\gamma_\varepsilon] = \int n(\gamma+\varepsilon\mathbf{h})|\mathbf{v}_\varepsilon|\,ds$ and differentiating with respect to $\varepsilon$, then setting $\varepsilon=0$, gives the first term $\int\nabla n\cdot\mathbf{h}\,ds$ and the second term $\int n\,\dot\gamma\cdot\dot{\mathbf{h}}\,ds$.
\end{proof}

\begin{theorem}[Ray equation]\label{thm:ray_equation}
Under Fermat's principle (\cref{def:fermat_principle}) and the slowly varying medium premise (\cref{rem:propagation_chapter_positioning}), the ray $\gamma(s)$ satisfies the \textbf{ray equation}
\[
\boxed{\ \frac{d}{ds}\Bigl(n(\gamma(s))\,\dot\gamma(s)\Bigr) = \nabla n(\gamma(s)).\ }
\]
\end{theorem}

\begin{proof}
By \cref{lem:optical_path_variation}, the first variation of the optical-path functional is
\[
\delta L[\mathbf{h}] = \int_{s_0}^{s_1}\nabla n\cdot\mathbf{h}\,ds + \int_{s_0}^{s_1}n\,\dot\gamma\cdot\dot{\mathbf{h}}\,ds.
\]
Integrate the second term by parts:
\[
\int_{s_0}^{s_1}n\,\dot\gamma\cdot\dot{\mathbf{h}}\,ds = \Bigl[n\,\dot\gamma\cdot\mathbf{h}\Bigr]_{s_0}^{s_1} - \int_{s_0}^{s_1}\frac{d}{ds}(n\dot\gamma)\cdot\mathbf{h}\,ds.
\]
By the fixed-endpoint conditions $\mathbf{h}(s_0)=\mathbf{h}(s_1)=\mathbf{0}$, the boundary term vanishes. Hence
\[
\delta L[\mathbf{h}] = \int_{s_0}^{s_1}\Bigl[\nabla n - \frac{d}{ds}(n\dot\gamma)\Bigr]\cdot\mathbf{h}\,ds.
\]
Fermat's principle requires that $\delta L[\mathbf{h}]=0$ for every $\mathbf{h}$ satisfying the endpoint conditions. By the fundamental lemma of the calculus of variations (\cref{lem:fundamental_lemma_calculus_of_variations}), the integrand vanishes identically:
\[
\nabla n - \frac{d}{ds}(n\dot\gamma) = \mathbf{0}. \qedhere
\]
\end{proof}

\begin{lemma}[Fundamental lemma of the calculus of variations]\label{lem:fundamental_lemma_calculus_of_variations}
Let $f\in C^0([s_0,s_1],\mathbb{R}^n)$. If $\int_{s_0}^{s_1}f(s)\cdot\mathbf{h}(s)\,ds = 0$ for every $C^1$ vector function $\mathbf{h}$ satisfying $\mathbf{h}(s_0)=\mathbf{h}(s_1)=\mathbf{0}$, then $f\equiv\mathbf{0}$.
\end{lemma}

\begin{proof}
Suppose to the contrary that there exists $s^*\in(s_0,s_1)$ with $f(s^*)\neq\mathbf{0}$. Take the $j$-th component $f_j(s^*)\neq 0$ of $f$. By continuity there exists $\delta>0$ such that $f_j(s)$ has the same sign as $f_j(s^*)$ whenever $|s-s^*|<\delta$. Take a cut-off function $\varphi\in C^1$ satisfying $\varphi(s)\ge 0$, $\varphi(s)=0$ for $|s-s^*|\ge\delta$ and $\varphi(s)>0$ for $|s-s^*|<\delta$. Put $\mathbf{h} = \varphi\,\mathbf{e}_j$ (where $\mathbf{e}_j$ is the $j$-th coordinate unit vector). Then $\int f\cdot\mathbf{h}\,ds = \int f_j\varphi\,ds$, and the integrand is everywhere non-negative on $|s-s^*|<\delta$ and is not identically zero, so the integral is strictly positive, contradicting the hypothesis. Hence $f\equiv\mathbf{0}$.
\end{proof}

\begin{corollary}[Eikonal equation]\label{cor:eikonal_equation}
If there exists an optical path function $S:\mathbb{R}^n\to\mathbb{R}$ such that the ray direction is $\dot\gamma = \nabla S/n$, then $S$ satisfies the \textbf{eikonal equation}
\[
\boxed{\ |\nabla S(\mathbf{r})|^2 = n(\mathbf{r})^2.\ }
\]
\end{corollary}

\begin{proof}
From $|\dot\gamma|=1$ and $\dot\gamma = \nabla S/n$, taking the square of the modulus gives $1 = |\nabla S|^2/n^2$, that is, $|\nabla S|^2 = n^2$.
\end{proof}

\begin{corollary}[Straight-line propagation in a homogeneous medium]\label{cor:straight_line_uniform}
If $n(\mathbf{r})\equiv n_0$ is constant, then the ray is a straight line and the propagation speed is $v = c/n_0$.
\end{corollary}

\begin{proof}
By \cref{thm:ray_equation}, $\nabla n = \mathbf{0}$, hence $d/ds(n_0\dot\gamma) = n_0\ddot\gamma = \mathbf{0}$, that is, $\ddot\gamma = \mathbf{0}$---the ray is a straight line. Propagation speed: from $ds/dt = c/n_0$ (which is given by the vacuum speed of light $c$ and the definition of the refractive index $n = c/v$), that is, $v = c/n_0$.
\end{proof}

\begin{proposition}[Wavelength scaling in a medium]\label{prop:medium_wavelength_scaling}
In a medium of refractive index $n$, the wavelength of light is scaled to $1/n$ of the vacuum wavelength:
\[
\lambda_{\text{medium}} = \frac{\lambda_{\text{vacuum}}}{n}.
\]
\end{proposition}

\begin{proof}
From the phase velocity of light in the medium, $v = c/n$, and the continuity of frequency across a medium interface (\cref{thm:snell_consequence}): $\lambda_{\text{medium}} = v/\nu = c/(n\nu) = \lambda_{\text{vacuum}}/n$.
\end{proof}

\begin{theorem}[Snell's law]\label{thm:snell_consequence}
At a flat interface between two homogeneous media, the angle of incidence $\theta_1$ and the angle of refraction $\theta_2$ satisfy
\[
n_1\sin\theta_1 = n_2\sin\theta_2.
\]
\end{theorem}

\begin{proof}
Apply Fermat's principle (\cref{def:fermat_principle}) to the optical path crossing the interface. Let the incident light travel from $\mathbf{A}$ to a point $\mathbf{P}$ on the interface, and then from $\mathbf{P}$ to $\mathbf{B}$. The optical path is $L(\mathbf{P}) = n_1|\mathbf{A}-\mathbf{P}| + n_2|\mathbf{P}-\mathbf{B}|$. Varying $\mathbf{P}$ within the interface and setting the first variation to zero gives $n_1\sin\theta_1 = n_2\sin\theta_2$ (where $\theta_1,\theta_2$ are respectively the angles between the optical path and the interface normal).
\end{proof}

\begin{figure}[htbp]
\centering
\begin{tikzpicture}[>=Stealth, thick, scale=1.0]
\fill[blue!8] (-3,0) rectangle (3,-2.2);
\fill[orange!8] (-3,0) rectangle (3,2.2);
\draw[thick] (-3,0) -- (3,0);
\node at (-2.3,0.25) {$n_1$};
\node at (-2.3,-0.35) {$n_2$};
\draw[dashed] (0,-2) -- (0,2);
\draw[->, blue, very thick] (-2,1.6) -- (0,0);
\draw[->, red, very thick] (0,0) -- (1.8,-1.4);
\draw[->, gray, dashed] (0,0) -- (2.2,0);
\node at (-0.25,0.9) {$\theta_1$};
\node at (0.35,-0.8) {$\theta_2$};
\node[blue] at (-1.5,1.35) {Incident};
\node[red] at (2.2,-1.1) {Refracted};
\node at (1.8,0.3) {Interface};
\node at (0.2,-2.0) {$n_1\sin\theta_1=n_2\sin\theta_2$};
\end{tikzpicture}
\caption{Snell's law and total internal reflection.}
\label{fig:snell}
\end{figure}

\begin{property}[Interface with the operators of the `Vector Analysis' chapter]\label{prop:geometric_optics_interface}
The basic equations of geometrical optics correspond exactly to the three operators of the `Vector Analysis' chapter:
\begin{enumerate}[label=(\roman*)]
    \item \textbf{Gradient}: the eikonal equation $|\nabla S|^2 = n^2$ (\cref{cor:eikonal_equation}) uses the gradient operator $\nabla$ directly. The ray direction is $\nabla S/n$, perpendicular to the wavefront $S=\text{const}$.
    \item \textbf{Divergence}: the flux conservation of a ray tube is given by the divergence theorem (\cref{thm:divergence}). Let $A(s)$ be the cross-sectional area of the ray tube; then $n(s)A(s)$ remains constant during propagation---this is the realisation in geometrical optics of the conservation laws of the `Vector Analysis' chapter.
    \item \textbf{Curl}: the amount by which a ray bends is given by the curl of the ray equation: $\nabla\times(n\dot\gamma) = \nabla\times(\nabla S) = \mathbf{0}$, that is, $n\dot\gamma$ is an irrotational field.
\end{enumerate}
\end{property}

\begin{proof}
(i) This has already been proved in \cref{cor:eikonal_equation}. (ii) Apply the divergence theorem to the integral over the lateral surface of the ray tube; combined with $\nabla\cdot(n\dot\gamma) = 0$ (computed directly from the ray equation), this gives $nA = \text{const}$. (iii) From $\dot\gamma = \nabla S/n$ we get $n\dot\gamma = \nabla S$, hence $\nabla\times(n\dot\gamma) = \nabla\times\nabla S = \mathbf{0}$ (\cref{thm:curl_grad_zero} of the `Vector Analysis' chapter).
\end{proof}

\begin{example}[Ray in a spherical-gradient medium]\label{ex:graded_index_ray}
Take $n(r) = n_0(1 - \alpha r^2)$ (a graded-index medium in the paraxial approximation). By the ray equation, a paraxial ray satisfies $\ddot{r} \approx -2\alpha n_0\,r$ (simple harmonic motion), and the ray follows a sinusoidal trajectory of period $T = 2\pi/\sqrt{2\alpha n_0}$. This is the periodic closure of rays in a self-focusing optical fibre---an instance in geometrical optics of `a finite-order action produces a closed orbit', consistent with the finite-order action framework of the `Operations and Periodicity' chapter.
\end{example}

\begin{figure}[htbp]
\centering
\begin{tikzpicture}[>=Stealth, thick, scale=0.9]
\draw[->] (-0.5,0) -- (7,0) node[right] {$z$};
\draw[->] (0,-1.5) -- (0,1.5) node[above] {$x$};
\draw[blue, very thick, domain=0:6.5, samples=200] plot (\x, {1.1*sin(deg(1.2*\x))});
\foreach \x in {0,1.57,3.14,4.71,6.28} {
  \draw[dashed, gray] (\x,-1.3) -- (\x,1.3);
}
\node at (3.2,-1.8) {Paraxial ray: $\ddot r \approx -2\alpha n_0 r$, period $T=2\pi/\sqrt{2\alpha n_0}$};
\end{tikzpicture}
\caption{Sinusoidal ray in a graded-index medium.}
\label{fig:graded_index_ray}
\end{figure}

\begin{remark}[Caustics and the failure of the eikonal equation]\label{rem:caustic_eikonal_failure}
The eikonal equation (\cref{cor:eikonal_equation}) fails on a \textbf{caustic surface}: a caustic surface is a surface on which the cross-sectional area of the ray tube satisfies $A(s)\to 0$, and on it the light intensity diverges as geometrical optics predicts ($I\propto 1/A\to\infty$). This divergence is not physical---the actual field strength is finite near the caustic and must be described by wave optics (the Airy function). Hence all geometrical-optics conclusions in this chapter take `away from the caustic surface' as a premise; the wave correction near a caustic is an \textbf{open problem} (\cref{rem:propagation_open}).
\end{remark}

\section{The propagation operator: an algebraic realisation of spatial translation}\label{sec:propagation_operator}

\begin{definition}[Free-space propagation operator]\label{def:propagation_operator}
Let $X$ be the function space in which the field quantities (scalar or vector components) live, taken to be $X = L^2(\mathbb{R}^n)$ or the Schwarz space $\mathcal{S}(\mathbb{R}^n)$. For $\mathbf{a}\in\mathbb{R}^n$, the \textbf{free-space propagation operator} is defined by
\[
(P_{\mathbf{a}} f)(\mathbf{x}) := f(\mathbf{x}-\mathbf{a}), \qquad f\in X,\ \mathbf{x}\in\mathbb{R}^n.
\]
\end{definition}

\begin{figure}[htbp]
\centering
\begin{tikzpicture}[>=Stealth, thick]
\draw[->] (-0.5,0) -- (6,0) node[right] {$\mathbf{x}$};
\draw[blue, very thick, domain=0:2.2, samples=100] plot (\x, {0.8*exp(-(\x-1)^2/0.4)});
\node[blue] at (1.1,1.1) {$f(\mathbf{x})$};
\draw[red, very thick, domain=2.5:4.7, samples=100] plot (\x, {0.8*exp(-(\x-3.5)^2/0.4)});
\node[red] at (3.5,1.1) {$(P_{\mathbf{a}}f)(\mathbf{x})=f(\mathbf{x}-\mathbf{a})$};
\draw[->, dashed] (1.0,-0.4) -- (3.4,-0.4) node[midway, below] {Translation $\mathbf{a}$};
\end{tikzpicture}
\caption{Illustration of the translation effected by the propagation operator.}
\label{fig:propagation_operator}
\end{figure}

\begin{definition}[Medium propagation operator]\label{def:medium_propagation_operator}
Let a homogeneous medium have refractive index $n_0$. Propagating a distance $\mathbf{a}$ in the medium (a geometrical distance) is equivalent to the optical distance $n_0\mathbf{a}$ in vacuum (the optical path). The \textbf{medium propagation operator} is defined by
\[
(P_{\mathbf{a}}^{(n_0)} f)(\mathbf{x}) := f(\mathbf{x}-n_0\mathbf{a}).
\]
\end{definition}

\begin{proposition}[Group structure of the propagation operator]\label{prop:propagation_group}
Under operator composition the family of propagation operators $\{P_{\mathbf{a}}\}_{\mathbf{a}\in\mathbb{R}^n}$ forms a faithful representation of the translation group $\mathbb{R}^n$:
\begin{enumerate}[label=(\roman*)]
    \item \textbf{Group homomorphism}: $P_{\mathbf{a}}P_{\mathbf{b}} = P_{\mathbf{a}+\mathbf{b}}$;
    \item \textbf{Identity element}: $P_{\mathbf{0}} = \operatorname{id}_X$;
    \item \textbf{Inverse element}: $P_{\mathbf{a}}^{-1} = P_{-\mathbf{a}}$.
\end{enumerate}
\end{proposition}

\begin{proof}
(i) $(P_{\mathbf{a}}P_{\mathbf{b}}f)(\mathbf{x}) = (P_{\mathbf{b}}f)(\mathbf{x}-\mathbf{a}) = f(\mathbf{x}-\mathbf{a}-\mathbf{b}) = (P_{\mathbf{a}+\mathbf{b}}f)(\mathbf{x})$. (ii)(iii) These follow from (i) on taking $\mathbf{b}=\mathbf{0}$ or $\mathbf{b}=-\mathbf{a}$.
\end{proof}

\begin{proposition}[Isometry of the propagation operator]\label{prop:propagation_isometry}
On $X = L^2(\mathbb{R}^n)$, the propagation operator $P_{\mathbf{a}}$ is an isometry:
\[
\|P_{\mathbf{a}} f\|_{L^2} = \|f\|_{L^2}, \qquad \forall f\in L^2(\mathbb{R}^n).
\]
In particular, $P_{\mathbf{a}}$ is a unitary operator.
\end{proposition}

\begin{proof}
Make the change of variable $\mathbf{y} = \mathbf{x}-\mathbf{a}$ (the Lebesgue measure is invariant):
\[
\|P_{\mathbf{a}} f\|_{L^2}^2 = \int_{\mathbb{R}^n}|f(\mathbf{x}-\mathbf{a})|^2\,d\mathbf{x} = \int_{\mathbb{R}^n}|f(\mathbf{y})|^2\,d\mathbf{y} = \|f\|_{L^2}^2.
\]
\end{proof}

\begin{proposition}[Fourier multiplier representation of the propagation operator]\label{prop:propagation_fourier}
Let $f\in L^2(\mathbb{R}^n)$ and let $\hat f(\mathbf{k}) = \int f(\mathbf{x})e^{-i\mathbf{k}\cdot\mathbf{x}}\,d\mathbf{x}$ be its Fourier transform. Then
\[
\widehat{P_{\mathbf{a}} f}(\mathbf{k}) = e^{-i\mathbf{k}\cdot\mathbf{a}}\hat f(\mathbf{k}).
\]
That is, in the Fourier domain the propagation operator is equivalent to multiplication by the phase factor $e^{-i\mathbf{k}\cdot\mathbf{a}}$.
\end{proposition}

\begin{proof}
Direct computation:
\[
\widehat{P_{\mathbf{a}} f}(\mathbf{k}) = \int f(\mathbf{x}-\mathbf{a})e^{-i\mathbf{k}\cdot\mathbf{x}}\,d\mathbf{x} = \int f(\mathbf{y})e^{-i\mathbf{k}\cdot(\mathbf{y}+\mathbf{a})}\,d\mathbf{y} = e^{-i\mathbf{k}\cdot\mathbf{a}}\hat f(\mathbf{k}). \qedhere
\]
\end{proof}

\begin{figure}[htbp]
\centering
\begin{tikzpicture}[>=Stealth, thick, node distance=1.8cm, scale=0.9, every node/.style={transform shape}]
\node[draw, rounded corners, fill=blue!5, align=center, minimum width=2.0cm] (x) {Spatial domain\\$f(\mathbf{x})$};
\node[draw, rounded corners, fill=orange!8, align=center, right=of x, minimum width=5.0cm] (p) {Propagation operator\\$P_{\mathbf{a}}f=f(\mathbf{x}-\mathbf{a})$};
\node[draw, rounded corners, fill=green!8, align=center, right=of p, minimum width=2.0cm] (k) {Frequency domain\\$\hat f(\mathbf{k})$};
\node[draw, rounded corners, fill=red!8, align=center, below=of p, minimum width=5.0cm] (phase) {Multiplication by a phase factor\\$e^{-i\mathbf{k}\cdot\mathbf{a}}\hat f(\mathbf{k})$};
\draw[->] (x) to[out=45, in=135] node[above] {$\mathcal{F}$} (k);
\draw[->] (x) -- (p);
\draw[->] (p) -- (phase);
\draw[->] (phase) -- (k);
\node at (5.3,-4.6) {When the propagation distance is $a=\lambda$, $\phi=2\pi$ and $e^{-i\phi}=1$: periodic closure.};
\end{tikzpicture}
\caption{The phase multiplier in the Fourier domain.}
\label{fig:fourier_phase}
\end{figure}

\begin{proposition}[Consistency of the propagation operator and the phase operator]\label{prop:phase_propagation}
For the vacuum plane wave $e^{i(\mathbf{k}\cdot\mathbf{x}-\omega t)}$, a propagation distance $a$ corresponds to the phase change $\phi = ka = \omega a/c$. In the frequency domain the propagation operator $P_a$ equals the phase operator
\[
\mathcal{O}_\phi^{(2\pi)} = e^{-i\phi},
\]
which has the same structure as the discrete phase operator of the `Periodic Functions and Groups with Periodic Orbits' chapter (\cref{def:discrete_phase_operator}).
\end{proposition}

\begin{proof}
By \cref{prop:propagation_fourier}, $P_a e^{ikx} = e^{-ika}e^{ikx}$. Taking $\phi = ka$ gives $e^{-i\phi}$.
\end{proof}

\begin{definition}[Complex exponential plane-wave representation of vacuum electromagnetic waves]\label{def:vacuum_em_wave}
In vacuum, the monochromatic solutions of the source-free Maxwell equations take the complex exponential plane-wave form
\[
\mathbf{E}(\mathbf{x},t)=\mathbf{E}_0\,e^{i(\mathbf{k}\cdot\mathbf{x}-\omega t)},\qquad
\mathbf{B}(\mathbf{x},t)=\mathbf{B}_0\,e^{i(\mathbf{k}\cdot\mathbf{x}-\omega t)},
\]
where the dispersion relation is $\omega=c|\mathbf{k}|$ and the transverse-wave conditions $\mathbf{k}\cdot\mathbf{E}_0=0$, $\mathbf{k}\times\mathbf{E}_0=\omega\mathbf{B}_0$ hold.
\end{definition}

\begin{corollary}[Periodic closure of propagation]\label{cor:propagation_period_closure}
When the propagation distance equals one wavelength, $a = \lambda = 2\pi/k$, the phase change is $\phi = 2\pi$ and $e^{-i\phi}=1$---the plane wave returns to the phase it had at departure. This is one complete \textbf{period} of the propagation operator.
\end{corollary}

\begin{proof}
By \cref{prop:phase_propagation}, $\phi = k\lambda = k\cdot 2\pi/k = 2\pi$ and $e^{-i2\pi}=1$.
\end{proof}

\begin{proposition}[The propagation operator commutes with constant-coefficient differential operators]\label{prop:propagation_commutes_translation}
Let $A$ be a linear differential operator with constant coefficients. Then
\[
P_{\mathbf{a}}A = A P_{\mathbf{a}}, \qquad \forall \mathbf{a}\in\mathbb{R}^n.
\]
\end{proposition}

\begin{proof}
It suffices to verify this for the monomials $A = \partial^\alpha$. By the chain rule, $\partial^\alpha(P_{\mathbf{a}}f) = P_{\mathbf{a}}(\partial^\alpha f)$. The result for polynomial sums follows by taking linear combinations.
\end{proof}

\begin{definition}[Evolution semigroup]\label{def:evolution_semigroup}
Let $X$ be a Banach space. A family of operators $\{E_t\}_{t\ge 0}\subset\mathcal{L}(X)$ is called a \textbf{strongly continuous evolution semigroup} if
\begin{enumerate}[label=(\roman*)]
    \item $E_0 = \operatorname{id}_X$;
    \item $E_{t+s} = E_t\circ E_s$ for all $t,s\ge 0$;
    \item for every $f\in X$, $\lim_{t\to 0^+}E_t f = f$ (strong continuity).
\end{enumerate}
Its \textbf{infinitesimal generator} $A$ is defined by $Af := \lim_{t\to 0^+}(E_t f - f)/t$, with domain consisting of all $f$ for which the above limit exists.
\end{definition}

\begin{figure}[htbp]
\centering
\begin{tikzpicture}[>=Stealth, thick]
\draw[->] (-0.5,0) -- (7,0) node[right] {$t$};
\foreach \t in {0,1,2,3,4,5,6} {
  \draw[gray!40] (\t,-0.15) -- (\t,0.15);
}
\node[blue] at (0.3,0.6) {$E_0=\operatorname{id}$};
\node[blue] at (2.0,0.6) {$E_s$};
\node[blue] at (4.0,0.6) {$E_t$};
\node[blue] at (6.0,0.6) {$E_{t+s}=E_t\circ E_s$};
\draw[->, red, very thick] (0,-0.7) -- (6,-0.7) node[midway, below] {Strongly continuous evolution semigroup};
\node at (3.0,-1.9) {Generator $A$: $\displaystyle Af=\lim_{t\to0^+}\frac{E_tf-f}{t}$, formally $E_t=e^{tA}$};
\end{tikzpicture}
\caption{Timeline of a strongly continuous evolution semigroup.}
\label{fig:evolution_semigroup}
\end{figure}

\begin{theorem}[Exponential representation of linear evolution equations]\label{thm:linear_evolution_exponential}
Let $X$ be a Banach space, let $A$ be a closed densely defined linear operator on $X$, and suppose that $\{E_t\}_{t\ge 0}$ is generated by $\partial_t u = Au$. Then formally
\[
E_t = e^{tA}, \qquad t\ge 0,
\]
holds on $\mathcal{D}(A)$.
\end{theorem}

\begin{proof}
For $u_0\in\mathcal{D}(A)$, the definition of the generator gives $E_t u_0 - u_0 = \int_0^t AE_s u_0\,ds$. Differentiating both sides with respect to $t$ gives $\frac{d}{dt}E_t u_0 = AE_t u_0$, that is, $E_t u_0$ is a solution of $\partial_t u = Au$. By the semigroup property, uniqueness of the solution guarantees that $E_t = e^{tA}$ (on $\mathcal{D}(A)$).
\end{proof}

\begin{theorem}[Energy conservation of vacuum light propagation]\label{thm:vacuum_energy_conservation}
Let $\mathbf{u}(t) = (\mathbf{E}(t),\mathbf{B}(t))$ be a solution of the vacuum Maxwell equations, and let the energy density integral be
\[
U(t) := \int_{\mathbb{R}^3}\Bigl(\frac{\varepsilon_0}{2}|\mathbf{E}(\mathbf{x},t)|^2 + \frac{1}{2\mu_0}|\mathbf{B}(\mathbf{x},t)|^2\Bigr)\,d\mathbf{x}.
\]
Then $U(t) = U(0)$ for all $t$.
\end{theorem}

\begin{proof}
Write the Maxwell equations as $\partial_t\mathbf{u} = A\mathbf{u}$, where $A$ is given by the curl operator. With respect to the $L^2$ energy inner product, $A$ is antisymmetric: $\langle A\mathbf{u},\mathbf{u}\rangle = 0$ (proved by integration by parts; this is Poynting's theorem, \cref{thm:electromagnetic_conservation} of the `Periodic Structure of Electromagnetic Systems' chapter). An antisymmetric operator on a Hilbert space generates a unitary group (Stone's theorem), hence $\|\mathbf{u}(t)\|^2 = \|\mathbf{u}(0)\|^2$, that is, $U(t)=U(0)$.
\end{proof}

\begin{corollary}[Mathematical formulation of causality]\label{cor:causality_speed}
Under the vacuum linear wave equation, the propagation speed of an electromagnetic signal does not exceed $c$. Concretely, if the initial field is supported at $t=0$ in an open set $\Omega_0$, then at $t>0$ the field is supported in the $ct$-neighbourhood $\Omega_0 + B_{ct}(\mathbf{0})$ of $\Omega_0$.
\end{corollary}

\begin{proof}
This follows from the unitary evolution group of \cref{thm:vacuum_energy_conservation} and the finite speed of propagation of the wave equation. Applying the divergence theorem to the energy flux shows that the energy crossing any sphere of radius $r$ centred at $\mathbf{x}_0$ is bounded by $c$ times the energy density outside the sphere; when the initial energy inside the sphere is zero, the energy inside the sphere remains zero for $t<r/c$.
\end{proof}

\begin{example}[Gaussian light pulse in vacuum]\label{ex:vacuum_gaussian_pulse}
Take the initial Gaussian pulse $\psi_0(x) = e^{-x^2/(2\sigma^2)}e^{ik_0x}$. The vacuum dispersion relation is $\omega(k) = ck$, the group velocity is $v_g = c$ and the second-order dispersion is $\omega''=0$. By \cref{thm:linear_evolution_exponential}, the pulse translates at speed $c$ without broadening:
\[
\psi(x,t) = e^{ik_0(x-ct)}e^{-(x-ct)^2/(2\sigma^2)}.
\]
This is the mathematical description of the phenomenon of `no dispersion'.
\end{example}

\begin{example}[Pulse broadening in a dispersive medium]\label{ex:dispersive_pulse}
In a dispersive medium the dispersion relation is $\omega(k)$ and the second-order dispersion is $\beta_2 = d^2\beta/d\omega^2$ (with $\beta = \omega n(\omega)/c$). A Gaussian pulse broadens:
\[
\sigma(t) \approx \sigma\sqrt{1 + (\beta_2 t/\sigma^2)^2}, \qquad t\to\infty.
\]
This is the dispersion penalty in optical fibre communication.
\end{example}

\begin{remark}[Envelope and carrier of a pulse]\label{rem:pulse_envelope}
The Gaussian pulses of \cref{ex:vacuum_gaussian_pulse} and \cref{ex:dispersive_pulse} are products of an envelope and a carrier: $e^{-(x-ct)^2/(2\sigma^2)}$ is the envelope and $e^{ik_0(x-ct)}$ the carrier. The group velocity $v_g=d\omega/dk$ is the propagation speed of the envelope, and the second-order dispersion $\beta_2$ describes its broadening---dispersion acts on the envelope rather than on the carrier (\cref{def:dyn_envelope}). This is the mathematical content of the slowly varying envelope approximation (SVEA): when the relative change of the envelope within one carrier period is much smaller than $1$, carrier and envelope can be treated separately.
\end{remark}

\section{The kernel-function system: integral representations of propagation}\label{sec:kernel_functions}

\begin{remark}[Positioning of this section]\label{rem:kernel_section_positioning}
This section systematically lists the kernel functions that occur throughout the book and classifies them by function. Kernel functions are the realisation of propagation and evolution at the integral level---every kernel function is the integral kernel of some class of propagation process or of some kind of spectral transform. This section distinguishes strictly between three classes of kernel function: (i) \textbf{propagation kernels} (Green's functions describing physical propagation); (ii) \textbf{window kernels and spectral kernels} (window functions and periodisation kernels in spectral analysis); (iii) \textbf{phase similarity kernels and weight kernels} (phase and frequency selection in periodic convolution).
\end{remark}

\begin{definition}[Propagation kernel (Green's function)]\label{def:propagation_kernel}
Suppose that the linear evolution equation $\partial_t u = \mathcal{L}u$ has a unique solution. The \textbf{propagation kernel} $G(\mathbf{x},t;\mathbf{x}_0)$ is the solution corresponding to the initial value $u(\mathbf{x},0) = \delta(\mathbf{x}-\mathbf{x}_0)$:
\[
u(\mathbf{x},t) = \int_{\mathbb{R}^n}G(\mathbf{x},t;\mathbf{x}_0)u_0(\mathbf{x}_0)\,d\mathbf{x}_0.
\]
Equivalently, $G$ is the integral kernel of the evolution operator $E_t = e^{t\mathcal{L}}$.
\end{definition}

\begin{theorem}[The propagation kernel satisfies the evolution equation and the initial condition]\label{thm:propagation_kernel_satisfies_pde}
The propagation kernel $G(\mathbf{x},t;\mathbf{x}_0)$ satisfies
\[
\partial_t G = \mathcal{L}G, \qquad G(\mathbf{x},0;\mathbf{x}_0) = \delta(\mathbf{x}-\mathbf{x}_0).
\]
\end{theorem}

\begin{proof}
By Definition \ref{def:propagation_kernel}, $u(\mathbf{x},t) = \int G(\mathbf{x},t;\mathbf{x}_0)u_0(\mathbf{x}_0)\,d\mathbf{x}_0$ satisfies $\partial_t u = \mathcal{L}u$ for every $u_0$. Taking $u_0 = \delta(\cdot-\mathbf{x}_0)$ gives the result. The initial condition is given by $E_0 = \operatorname{id}$.
\end{proof}

\begin{longtable}{@{}>{\raggedright\arraybackslash\hspace{0pt}}p{2.3cm}
                >{\raggedright\arraybackslash\hspace{0pt}\small}p{5.2cm}
                >{\raggedright\arraybackslash\hspace{0pt}}p{3.8cm}
                >{\raggedright\arraybackslash\hspace{0pt}}p{2.2cm}@{}}
\caption{Complete table of the kernel functions of the book (classified by function)}
\label{tab:all_kernels}\\
\toprule
Kernel function & Mathematical form & Physical/mathematical meaning & Chapter \\
\midrule
\endfirsthead

\multicolumn{4}{@{}l}{\small\itshape(continued)}\\
\toprule
Kernel function & Mathematical form & Physical/mathematical meaning & Chapter \\
\midrule
\endhead

\midrule
\multicolumn{4}{r@{}}{\small\itshape continued on next page}\\
\endfoot

\bottomrule
\endlastfoot

\multicolumn{4}{@{}l}{\textbf{(A) Propagation kernels (Green's functions)}} \\
\midrule
Free-space Green's function & $G(\mathbf{x},\mathbf{x}') = \frac{e^{ik\|\mathbf{x}-\mathbf{x}'\|}}{4\pi\|\mathbf{x}-\mathbf{x}'\|}$ & Helmholtz propagation of a point source in vacuum & Electromagnetic systems \\
Heat kernel & $G(x,t) = \frac{1}{\sqrt{4\pi\alpha t}}e^{-x^2/(4\alpha t)}$ & Diffusion kernel of heat conduction & Differential equations \\
One-dimensional wave propagation kernel & $G(x,t) = \frac{1}{2c}\mathbf{1}_{|x|<ct}$ & Fundamental solution of the one-dimensional wave equation & Differential equations \\
Poisson kernel & $P_r(\theta) = \frac{1-r^2}{1-2r\cos\theta+r^2}$ & Propagation of boundary values of harmonic functions & Harmonic analysis \\
Retarded Green's function & $G_{\text{ret}} = \frac{\delta(t'-t+\|\mathbf{x}-\mathbf{x}'\|/c)}{4\pi\|\mathbf{x}-\mathbf{x}'\|}$ & Integral kernel of the retarded potential & Electromagnetic systems \\
Airy kernel & $\operatorname{Ai}(x) = \frac{1}{\pi}\int_0^\infty\cos(t^3/3+xt)\,dt$ & Wave correction near a caustic & Note to this chapter \\
\midrule
\multicolumn{4}{@{}l}{\textbf{(B) Window kernels and spectral kernels}} \\
\midrule
Rectangular window kernel & $W_T(\omega) = \frac{1}{T}\frac{\sin(T\omega/2)}{\sin(\omega/2)}$ & Spectral kernel of a finite window, sidelobes $-13$ dB & Discrete sampling \\
Hann window kernel & $w(n) = \tfrac12\bigl(1-\cos\tfrac{2\pi n}{M-1}\bigr)$ & Sidelobes $-31$ dB, decay $1/\omega^3$ & Discrete sampling \\
Hamming window kernel & $w(n) = 0.54-0.46\cos\tfrac{2\pi n}{M-1}$ & Sidelobes $-43$ dB & Discrete sampling \\
Blackman window kernel & $w(n) = 0.42-0.5\cos\tfrac{2\pi n}{M-1}+0.08\cos\tfrac{4\pi n}{M-1}$ & Sidelobes $-58$ dB & Discrete sampling \\
Triangular window kernel & $w(n) = 1-\bigl|\tfrac{2n}{M-1}-1\bigr|$ & Sidelobes $-27$ dB, decay $1/\omega^2$ & Discrete sampling \\
Gaussian kernel & $K_\Delta(x) = \frac{1}{\sqrt{2\pi}\Delta}e^{-x^2/(2\Delta^2)}$ & Optimal in time and frequency, no sidelobes & Time series analysis \\
Mexican-hat kernel & $K_{\text{MexHat}}(x) = (1-x^2/\Delta^2)e^{-x^2/(2\Delta^2)}$ & Band-pass, second derivative & Time series analysis \\
Cauchy kernel & $K_\Delta(x) = \frac{1}{\pi}\frac{\Delta}{x^2+\Delta^2}$ & Heavy-tailed, long-range & Time series analysis \\
Sinc kernel & $K(x) = \frac{\sin(\pi x/\Delta)}{\pi x/\Delta}$ & Ideal low-pass & Harmonic analysis \\
Dirichlet kernel & $D_M(\theta) = \frac{\sin((M+\frac12)\theta)}{\sin(\theta/2)}$ & Periodic window kernel, main lobe $2\pi/T$ & Discrete sampling \\
Fejér kernel & $F_N(x) = \frac{1}{N+1}\Bigl(\frac{\sin((N+1)x/2)}{\sin(x/2)}\Bigr)^2$ & Non-negative good kernel, Cesàro means & Harmonic analysis \\
\midrule
\multicolumn{4}{@{}l}{\textbf{(C) Phase similarity kernels (in periodic convolution)}} \\
\midrule
$S^{(1)}$ & $e^{i(\theta_m-\varphi_n)}$ & Basic difference angle (modulus of the inner/outer product) & Scalars and vectors \\
$S^{(2)}$ & $e^{i(\theta_m+\varphi_n)}$ & Basic sum angle (complex multiplication) & Scalars and vectors \\
$S^{(3)}$ & $e^{i2(\theta_m-\varphi_n)}$ & Double-angle difference angle (second-order reversal detection) & Scalars and vectors \\
$S^{(4)}$ & $e^{i2(\theta_m+\varphi_n)}$ & Double-angle sum angle & Scalars and vectors \\
$S^{(5)}$ & $e^{i(\theta_m-\varphi_n)/2}$ & Half-angle difference angle (modulus of the sum/difference vector) & Scalars and vectors \\
$S^{(6)}$ & $e^{i(\theta_m+\varphi_n)/2}$ & Half-angle sum angle & Scalars and vectors \\
\midrule
\multicolumn{4}{@{}l}{\textbf{(D) Weight kernels and periodic kernels}} \\
\midrule
Weight kernel $W$ & $W(\omega_m-\nu_n)$ & Frequency-domain selection (bandwidth control) & Series and convolution \\
Periodic kernel $k_{\text{per}}$ & $\sigma^2\exp\bigl(-\tfrac{2\sin^2(\pi|x-x'|/T)}{\ell^2}\bigr)$ & GP periodic prior (MacKay) & Probabilistic systems \\
Gaussian-process kernel $k_{\text{SE}}$ & $\sigma^2 e^{-(x-x')^2/(2\ell^2)}$ & GP smooth prior & Probabilistic systems \\
\end{longtable}

\begin{figure}[htbp]
\centering
\begin{fittikz}
\begin{tikzpicture}[
  every node/.style={draw, rounded corners, align=center, fill=blue!4,
                     font=\footnotesize, minimum width=2.2cm, minimum height=0.8cm},
  edge/.style={draw, -{Stealth[length=2mm]}, thick}
]
\node (root) at (0,0) {Kernel-function system};

\node (TL) at (-4.1,3)  {Propagation kernels\\Green's functions};

\node (TR) at (3.5,2)   {Window and spectral kernels};
% The phase similarity kernels, the weight kernels and the periodic kernels are shifted 2cm to the left as a whole
\node (BL) at (-6.0,-2) {Phase similarity kernels};
\node (BR) at (2.5,-2)  {Weight and periodic kernels};

\draw[edge] (root) -- (TL);
\draw[edge] (root) -- (TR);
\draw[edge] (root) -- (BL);
\draw[edge] (root) -- (BR);

\node (TL1) at (-9.0,5.0) {Free-space\\Green's function};
\node (TL2) at (-9.0,3.8) {Heat kernel};
\node (TL3) at (-9.0,2.6) {Wave propagation kernel};
\node (TL4) at (-9.0,1.4) {Poisson kernel};
\node (TL5) at (-9.0,0.2) {Retarded Green's function};
\draw[edge] (TL) -- (TL1);
\draw[edge] (TL) -- (TL2);
\draw[edge] (TL) -- (TL3);
\draw[edge] (TL) -- (TL4);
\draw[edge] (TL) -- (TL5);

\node (TR1) at (0.0,4) {Rectangular/Hann\\Hamming/Blackman};
\node (TR2) at (3.5,4) {Gaussian/Mexican hat\\Cauchy/sinc};
\node (TR3) at (3.5,0) {Dirichlet\\Fejér};
\draw[edge] (TR) -- (TR1);
\draw[edge] (TR) -- (TR2);
\draw[edge] (TR) -- (TR3);

\node (BL1) at (-8.5,-4) {$S^{(1)},S^{(2)}$};
\node (BL2) at (-6.0,-4) {$S^{(3)},S^{(4)}$};
\node (BL3) at (-3.5,-4) {$S^{(5)},S^{(6)}$};
\draw[edge] (BL) -- (BL1);
\draw[edge] (BL) -- (BL2);
\draw[edge] (BL) -- (BL3);

\node (BR1) at (0.0,-4) {Weight kernel $W$};
\node (BR2) at (2.5,-4) {Periodic kernel $k_{\text{per}}$};
\node (BR3) at (5.0,-4) {GP Gaussian kernel};
\draw[edge] (BR) -- (BR1);
\draw[edge] (BR) -- (BR2);
\draw[edge] (BR) -- (BR3);

\end{tikzpicture}
\end{fittikz}
\caption{Classification tree of kernel functions.}
\label{fig:kernel_tree}
\end{figure}

\begin{property}[Duality relations among kernel functions]\label{prop:kernel_duality}
The kernel functions satisfy the following duality relations:
\begin{enumerate}[label=(\roman*)]
    \item \textbf{Gaussian--Gaussian duality}: the Fourier transform of a Gaussian window kernel is again a Gaussian kernel (Lemma \ref{lem:gaussian_fourier_duality}).
    \item \textbf{Window--kernel duality}: the rectangular window kernel and the sinc kernel are Fourier duals of each other (Theorem \ref{thm:window_kernel_duality}).
    \item \textbf{Phase-kernel conjugacy}: $S^{(d)}$ and $S^{(d')}$ satisfy the conjugacy duality $S^{(1)} = \overline{S^{(1)}_{\text{reverse}}}$ (\cref{prop:phase_convolution_symmetry}).
    \item \textbf{Propagation--evolution duality}: the spatial propagation kernel and the temporal evolution kernel are spacetime duals of each other (\cref{thm:propagation_evolution_duality}).
\end{enumerate}
\end{property}

\begin{proof}
(i) Compute the Fourier transform of the Gaussian kernel directly. (ii) By the window--kernel duality theorem of the `Discrete Sampling and Intercalation of Periodic Signals' chapter (\cref{thm:window_kernel_duality}). (iii) By the conjugate symmetry of phase convolution in the `Series and Convolution' chapter (\cref{prop:phase_convolution_symmetry}). (iv) See \cref{thm:propagation_evolution_duality}.
\end{proof}

\begin{lemma}[Fourier self-duality of the Gaussian kernel]\label{lem:gaussian_fourier_duality}
Let $K_\Delta(x) = \frac{1}{\sqrt{2\pi}\Delta}e^{-x^2/(2\Delta^2)}$. Then
\[
\hat K_\Delta(\omega) = e^{-\Delta^2\omega^2/2},
\]
that is, the Fourier transform of a Gaussian kernel is again a Gaussian kernel (the variances are interchanged, $\Delta\leftrightarrow 1/\Delta$).
\end{lemma}

\begin{proof}
By the standard Gaussian integral $\int_{-\infty}^\infty e^{-x^2/(2\Delta^2)}e^{-i\omega x}\,dx = \sqrt{2\pi}\Delta e^{-\Delta^2\omega^2/2}$ (complete the square and then perform the Gaussian integral).
\end{proof}

\begin{theorem}[Spacetime duality of the propagation kernel and the evolution kernel]\label{thm:propagation_evolution_duality}
Let $G_{\text{prop}}(\mathbf{x})$ be the spatial propagation kernel (the Green's function of the Helmholtz equation) and let $G_{\text{evo}}(t)$ be the temporal evolution kernel. In the Fourier domain the two are linked by the spacetime dispersion relation $\omega(\mathbf{k})$:
\[
\widehat{G_{\text{prop}}}(\mathbf{k}) = \frac{1}{\omega(\mathbf{k})}, \qquad \widehat{G_{\text{evo}}}(\omega) = \frac{1}{\omega},
\]
and the spacetime convolution of $G_{\text{prop}}$ with $G_{\text{evo}}$ gives the complete propagation--evolution kernel.
\end{theorem}

\begin{proof}
In the Fourier domain the Helmholtz equation $(\nabla^2+k^2)G = -\delta$ becomes $(-|\mathbf{k}|^2+k^2)\hat G = -1$, that is, $\hat G(\mathbf{k}) = 1/(|\mathbf{k}|^2-k^2)$. For the temporal evolution equation $\partial_t u = i\omega u$, the Fourier transform gives $\hat G_{\text{evo}}(\omega) = 1/(-i\omega)$, whose modulus is $1/|\omega|$. The spacetime duality is provided by the dispersion relation $\omega = \omega(\mathbf{k})$.
\end{proof}

\begin{example}[Far-field expansion of the free-space Green's function]\label{ex:green_far_field}
In the far field $|\mathbf{x}|\gg|\mathbf{x}'|$ the free-space Green's function $G(\mathbf{x},\mathbf{x}') = e^{ik\|\mathbf{x}-\mathbf{x}'\|}/(4\pi\|\mathbf{x}-\mathbf{x}'\|)$ expands as
\[
G(\mathbf{x},\mathbf{x}') \approx \frac{e^{ik|\mathbf{x}|}}{4\pi|\mathbf{x}|}e^{-i\mathbf{k}\cdot\mathbf{x}'}, \qquad \mathbf{k} = k\hat{\mathbf{x}}.
\]
This is the physical origin of the phase factor $e^{-i\mathbf{k}\cdot\mathbf{a}}$---the realisation, in the far-field expansion of the Green's function, of the Fourier multiplier representation of the propagation operator (\cref{prop:propagation_fourier}).
\end{example}

\section{Propagation convolution and evolution convolution}\label{sec:propagation_convolution}

\begin{remark}[Positioning of this section]\label{rem:convolution_section_positioning}
Building on the kernel-function system (Section \ref{sec:kernel_functions}), this section defines \textbf{propagation convolution} and \textbf{evolution convolution}. Their common structure is that the action of a propagation operator or of an evolution operator on a function is equivalent to the convolution of the function with the corresponding kernel. The strict premises of this section agree with the rest of the chapter: a linear evolution equation plus a slowly varying medium. At the end of the section the relation of the two to the periodic convolution of the `Series and Convolution' chapter is made explicit.
\end{remark}

\begin{definition}[Propagation convolution]\label{def:propagation_convolution}
Let $G_{\text{prop}}(\mathbf{x},\mathbf{x}')$ be a spatial propagation kernel (\cref{def:propagation_kernel}). The \textbf{propagation convolution} of a function $f$ with $G$ is defined by
\[
(f \circledast_{\text{prop}} G)(\mathbf{x}) := \int_{\mathbb{R}^n}f(\mathbf{x}')G(\mathbf{x},\mathbf{x}')\,d\mathbf{x}'.
\]
When the medium is homogeneous (translation invariant), $G(\mathbf{x},\mathbf{x}') = G(\mathbf{x}-\mathbf{x}')$ and propagation convolution degenerates to ordinary convolution
\[
(f \circledast_{\text{prop}} G)(\mathbf{x}) = \int_{\mathbb{R}^n}f(\mathbf{x}')G(\mathbf{x}-\mathbf{x}')\,d\mathbf{x}' = (f * G)(\mathbf{x}).
\]
\end{definition}

\begin{definition}[Evolution convolution]\label{def:evolution_convolution}
Let $G_{\text{evo}}(t,t')$ be a temporal evolution kernel (the temporal dual of \cref{def:propagation_kernel}). The \textbf{evolution convolution} of a function $f$ with $G$ is defined by
\[
(f \circledast_{\text{evo}} G)(t) := \int_{\mathbb{R}}f(t')G(t,t')\,dt'.
\]
When the evolution equation has constant coefficients (is invariant under time translation), $G(t,t') = G(t-t')$ and evolution convolution degenerates to ordinary time convolution
\[
(f \circledast_{\text{evo}} G)(t) = (f * G)(t).
\]
\end{definition}

\begin{definition}[Joint propagation--evolution convolution]\label{def:joint_convolution}
Let $G(\mathbf{x},t;\mathbf{x}',t')$ be a spacetime propagation kernel. The \textbf{joint propagation--evolution convolution} of a function $f$ with $G$ is defined by
\[
(f \circledast_{\text{PE}} G)(\mathbf{x},t) := \int_{\mathbb{R}^n}\int_{\mathbb{R}}f(\mathbf{x}',t')G(\mathbf{x},t;\mathbf{x}',t')\,d\mathbf{x}'\,dt'.
\]
When the medium is homogeneous and the evolution has constant coefficients, $G(\mathbf{x},t;\mathbf{x}',t') = G(\mathbf{x}-\mathbf{x}', t-t')$ and the joint convolution is the ordinary spacetime convolution.
\end{definition}

\begin{theorem}[Relation between propagation convolution and periodic convolution]\label{thm:propagation_periodic_convolution}
The periodic convolution of the single-period layer ($P=Q=\{T\}$, $W\to\delta$) (\cref{def:periodic_convolution}) degenerates on the spatial translation group $\mathbb{R}^n$ to propagation convolution:
\[
(f \circledast_{\text{prop}} G)(\mathbf{x}) = \lim_{W\to\delta}(f\circledast g)(\mathbf{x}),
\]
where $g$ is taken to be the propagation kernel $G$ and $S^{(d)}$ is taken to be the phase kernel $S^{(1)}$.
\end{theorem}

\begin{proof}
By the single-period-layer degeneration of \cref{def:periodic_convolution} (\cref{prop:periodic_convolution_degenerations}(i)) and \cref{def:propagation_convolution}, both are convolutions over the same frequency pair $(n=m)$. The phase kernel satisfies $S^{(1)}(\theta,\theta)=1$ and the amplitude switch $\alpha=1$ preserves the amplitude. Hence periodic convolution degenerates to propagation convolution.
\end{proof}

\begin{figure}[htbp]
\centering
\begin{fittikz}
\begin{tikzpicture}[>=Stealth, thick, node distance=1.2cm]
\node[draw, rounded corners, fill=purple!8, align=center] (pc) {Periodic convolution\\$f\circledast g$};
\node[draw, rounded corners, fill=blue!8, align=center, below left=of pc] (prop) {Propagation convolution\\$f\circledast_{\text{prop}}G$};
\node[draw, rounded corners, fill=green!8, align=center, below right=of pc] (evo) {Evolution convolution\\$f\circledast_{\text{evo}}G$};
\node[draw, rounded corners, fill=orange!8, align=center, below=of pc] (joint) {Joint propagation--evolution convolution\\$f\circledast_{\text{PE}}G$};
\draw[->] (pc) -- node[left, align=center] {$W\to\delta$\\$S^{(1)}=1$\\$\alpha=1$} (prop);
\draw[->] (pc) -- node[right, align=center] {Time translation\\single-period degeneration} (evo);
\draw[->] (prop) -- (joint);
\draw[->] (evo) -- (joint);
\node at (0,-4.2) {In the single-period degeneration the phase kernel, the weight kernel and the amplitude switch of periodic convolution become $S^{(1)}=1$, $W=\delta$ and $\alpha=1$ respectively.};
\end{tikzpicture}
\end{fittikz}
\caption{Degeneration from periodic convolution to propagation convolution.}
\label{fig:convolution_degeneration}
\end{figure}

\begin{property}[Algebraic properties of propagation convolution]\label{prop:propagation_convolution_algebra}
In a translation-invariant medium, propagation convolution satisfies:
\begin{enumerate}[label=(\roman*)]
    \item \textbf{Commutativity}: $f\circledast_{\text{prop}}G = G\circledast_{\text{prop}}f$ (by the commutativity of convolution);
    \item \textbf{Associativity}: $(f\circledast_{\text{prop}}G_1)\circledast_{\text{prop}}G_2 = f\circledast_{\text{prop}}(G_1\circledast_{\text{prop}}G_2)$;
    \item \textbf{Fourier multiplier}: $\widehat{f\circledast_{\text{prop}}G}(\mathbf{k}) = \hat f(\mathbf{k})\hat G(\mathbf{k})$;
    \item \textbf{Realisation by the propagation operator}: $P_{\mathbf{a}}f = f\circledast_{\text{prop}}\delta_{\mathbf{a}}$, where $\delta_{\mathbf{a}}(\mathbf{x}) = \delta(\mathbf{x}-\mathbf{a})$.
\end{enumerate}
\end{property}

\begin{proof}
(i)(ii) By the commutativity and associativity of ordinary convolution. (iii) By the convolution theorem. (iv) By the sifting property of the $\delta$ function: $\int f(\mathbf{x}')\delta(\mathbf{x}-\mathbf{a}-\mathbf{x}')\,d\mathbf{x}' = f(\mathbf{x}-\mathbf{a}) = P_{\mathbf{a}}f(\mathbf{x})$.
\end{proof}

\begin{theorem}[Unification of propagation convolution and the propagation operator]\label{thm:propagation_convolution_operator}
Let $G_{\text{prop}}$ be the propagation kernel in a homogeneous medium. Then the action of the propagation operator $P_{\mathbf{a}}$ on $f$ is equivalent to the propagation convolution of $f$ with $\delta_{\mathbf{a}}$:
\[
P_{\mathbf{a}}f = f\circledast_{\text{prop}}\delta_{\mathbf{a}}.
\]
More generally, if $G$ is an arbitrary propagation kernel, then $f\circledast_{\text{prop}}G$ gives the propagation result with $G$ as propagation kernel.
\end{theorem}

\begin{proof}
The first formula has already been proved in \cref{prop:propagation_convolution_algebra}(iv). The second follows directly from Definition \ref{def:propagation_convolution} and the definition of the propagation kernel, Definition \ref{def:propagation_kernel}.
\end{proof}

\begin{example}[Evolution convolution of heat conduction]\label{ex:heat_evolution_convolution}
The propagation kernel of the one-dimensional heat equation $\partial_t u = \alpha\partial_x^2 u$ is the heat kernel $G(x,t) = \frac{1}{\sqrt{4\pi\alpha t}}e^{-x^2/(4\alpha t)}$. The evolution convolution
\[
u(x,t) = (u_0\circledast_{\text{evo}}G)(x,t) = \int_{\mathbb{R}}u_0(x')\frac{1}{\sqrt{4\pi\alpha t}}e^{-(x-x')^2/(4\alpha t)}\,dx'
\]
gives the evolution of the initial value $u_0$. By kernel duality (\cref{prop:kernel_duality}), the Fourier transform of the heat kernel is $e^{-\alpha k^2 t}$---this is the exponential decay factor of the frequency components, consistent with the Fourier solution of the heat equation in the `Differential Equations and Analysis' chapter (\cref{ex:heat_fourier}).
\end{example}

\begin{example}[Propagation convolution of free-space light]\label{ex:free_space_optical_convolution}
The propagation kernel of light in vacuum is the free-space Green's function $G(\mathbf{x}) = e^{ik|\mathbf{x}|}/(4\pi|\mathbf{x}|)$. The propagation convolution
\[
u(\mathbf{x}) = (u_0\circledast_{\text{prop}}G)(\mathbf{x}) = \int_{\mathbb{R}^3}u_0(\mathbf{x}')\frac{e^{ik|\mathbf{x}-\mathbf{x}'|}}{4\pi|\mathbf{x}-\mathbf{x}'|}\,d\mathbf{x}'
\]
gives the solution of the Helmholtz equation---this is the mathematical form of Kirchhoff's diffraction formula.
\end{example}

\begin{remark}[Summary of the relation between propagation convolution and periodic convolution]\label{rem:propagation_convolution_summary}
The relation between propagation convolution and periodic convolution may be summarised as follows: \textbf{propagation convolution is the single-period degeneration of periodic convolution on the spatial translation group} (\cref{thm:propagation_periodic_convolution}), and \textbf{evolution convolution is the single-period degeneration of periodic convolution on the time translation group}. The combination of the two (\cref{def:joint_convolution}) gives complete spacetime propagation. In propagation convolution the phase similarity kernel $S^{(d)}$, the weight kernel $W$ and the amplitude switch $\alpha$ of periodic convolution degenerate respectively to: the phase kernel $S^{(1)}=1$, the weight kernel $W=\delta$ and $\alpha=1$.
\end{remark}

\begin{remark}[Fourier representation of the joint propagation--evolution convolution]\label{rem:joint_convolution_fourier}
By the product theorem for Fourier transforms, the joint propagation--evolution convolution in the spacetime frequency domain is
\[
\widehat{f\circledast_{\text{PE}}G}(\mathbf{k},\omega) = \hat f(\mathbf{k},\omega)\hat G(\mathbf{k},\omega).
\]
Here $\hat G(\mathbf{k},\omega)$ is determined by the spacetime dispersion relation $\omega = \omega(\mathbf{k})$. In particular, the free-space $\hat G$ takes non-zero values on $\omega = c|\mathbf{k}|$---this is precisely a restatement, in the framework of joint convolution, of the plane-wave dispersion relation of the `Periodic Structure of Electromagnetic Systems' chapter (\cref{prop:plane_wave_periodic}).
\end{remark}

\begin{theorem}[Bloch--Floquet propagation]\label{thm:bloch_floquet}
Let the linear evolution operator $\mathcal{L}$ have spatial period $\mathbf{T}$ (that is, $P_{\mathbf{T}}\mathcal{L} = \mathcal{L}P_{\mathbf{T}}$). Then there exist Bloch wave solutions of the form
\[
u_{\mathbf{k},n}(\mathbf{x},t) = e^{i(\mathbf{k}\cdot\mathbf{x}-\omega_n(\mathbf{k})t)}\phi_{\mathbf{k},n}(\mathbf{x}), \qquad \phi_{\mathbf{k},n}(\mathbf{x}+\mathbf{T}) = \phi_{\mathbf{k},n}(\mathbf{x})
\]
Each $n$ corresponds to an \textbf{energy band}, $\omega_n(\mathbf{k})$ is the dispersion relation of the $n$-th energy band, and $\mathbf{k}$ ranges over the \textbf{first Brillouin zone} $\mathbb{R}^n/\mathbf{T}^*\mathbb{Z}^n$.
\end{theorem}

\begin{proof}
Since $\mathcal{L}$ commutes with $P_{\mathbf{T}}$, the two admit a simultaneous spectral decomposition on $L^2$. The operator $P_{\mathbf{T}}$ is unitary and its spectrum is the unit circle; each eigenvalue may be written in the form $e^{-i\mathbf{k}\cdot\mathbf{T}}$ (with $\mathbf{k}$ taken modulo $\mathbf{T}^*$). The corresponding eigenspace is spanned by Bloch waves: $u(\mathbf{x}) = e^{i\mathbf{k}\cdot\mathbf{x}}\phi(\mathbf{x})$, where $\phi$ is periodic with period $\mathbf{T}$. Grading further by $\omega_n(\mathbf{k})$ in the spectral decomposition of $\mathcal{L}$ gives the result.
\end{proof}

\begin{figure}[htbp]
\centering
\begin{tikzpicture}[>=Stealth, thick, scale=0.95]
\draw[->] (-0.5,0) -- (6.5,0) node[right] {$k$};
\draw[->] (0,-0.2) -- (0,3.2) node[above] {$\omega$};
\foreach \x in {0,1,2,3,4,5,6} {
  \draw[gray!30] (\x,0) -- (\x,3.0);
}
\draw[blue, very thick, domain=0:3, samples=100] plot (\x, {0.5+0.4*sin(deg(1.2*\x))});
\draw[red, very thick, domain=0:3, samples=100] plot (\x, {1.6+0.35*sin(deg(1.5*\x+0.5))});
\draw[green!60!black, very thick, domain=0:3, samples=100] plot (\x, {2.5+0.3*sin(deg(1.0*\x+1.0))});
\node[blue] at (1.5,0.2) {Energy band $n=1$};
\node[red] at (1.5,1.3) {Energy band $n=2$};
\node[green!60!black] at (1.5,2.2) {Energy band $n=3$};
\draw[dashed] (0,0.95) -- (3,0.95);
\draw[dashed] (0,2.05) -- (3,2.05);
\node at (4.5,1.5) {Band gap};
\draw[<->] (3.2,0.95) -- (3.2,2.05);
\node at (3.0,-0.7) {First Brillouin zone $\mathbb{R}^n/\mathbf{T}^*\mathbb{Z}^n$};
\end{tikzpicture}
\caption{Bloch energy bands and the Brillouin zone.}
\label{fig:bloch_bands}
\end{figure}

\begin{corollary}[Photonic band gap]\label{cor:photonic_bandgap}
In a photonic crystal whose dielectric constant is periodic with period $\mathbf{T}$, if some frequency interval contains no energy band for any real $\mathbf{k}$, then that interval is a \textbf{photonic band gap}. Light in the band gap is forbidden to propagate in the crystal.
\end{corollary}

\begin{proof}
By \cref{thm:bloch_floquet}, the light propagation modes are described by the energy bands $\omega_n(\mathbf{k})$. If a frequency interval does not lie in the range of any $\omega_n(\mathbb{R}^n)$, then no corresponding Bloch wave solution exists and light cannot exist as a propagating mode.
\end{proof}

\begin{property}[Frequency-domain support of the dispersion relation]\label{prop:dispersion_period}
In a periodic medium the dispersion relation $\omega_n(\mathbf{k})$ is a periodic function on the first Brillouin zone:
\[
\omega_n(\mathbf{k}+\mathbf{G}) = \omega_n(\mathbf{k}), \qquad \forall \mathbf{G}\in\mathbf{T}^*\mathbb{Z}^n.
\]
That is, the band structure has period $\mathbf{G}$ in the frequency domain---this is the representation, in the frequency domain, of the spatial period $\mathbf{T}$ (Pontryagin duality, \cref{thm:three_groups_duality} of the `Groups and Group Functions' chapter).
\end{property}

\begin{proof}
By the definition of the Bloch wave, $u_{\mathbf{k},n} = e^{i(\mathbf{k}\cdot\mathbf{x}-\omega_n(\mathbf{k})t)}\phi_{\mathbf{k},n}(\mathbf{x})$; when $\mathbf{k}\to\mathbf{k}+\mathbf{G}$ (with $\mathbf{G}$ a reciprocal lattice vector), the product of $e^{i\mathbf{G}\cdot\mathbf{x}}$ with $\phi_{\mathbf{k},n}$ is still a function of period $\mathbf{T}$, hence $\omega_n(\mathbf{k}+\mathbf{G}) = \omega_n(\mathbf{k})$.
\end{proof}

\begin{remark}[Duality with Floquet theory of the `Differential Equations and Analysis' chapter]\label{rem:propagation_floquet}
\cref{thm:bloch_floquet} is the dual, in a spatially periodic medium, of Floquet's theorem (\cref{thm:floquet}) in the `Differential Equations and Analysis' chapter: the Floquet multiplier of a time-periodic system $\leftrightarrow$ the Bloch phase of a spatially periodic system; the time evolution group $\leftrightarrow$ the spatial propagation operator. Both are spectral theories of linear systems with periodic coefficients; only the direction of the periodicity differs.
\end{remark}

\begin{remark}[Interface with the `Periodic Manifold' chapter]\label{rem:propagation_periodic_manifold}
In the `Periodic Manifold' chapter, a discrete unitary evolution $\mathcal{U}^p=I$ on lattice points gives a strict period. The evolution semigroup of this chapter is its continuum limit: as the time step $\Delta t\to 0$, the discrete evolution tends to the continuous evolution group $E_t = e^{tA}$ (\cref{thm:continuum_limit_full}). Bloch waves correspond on a periodic manifold to periodic orbits on lattice points---they already appear in the continuum limit of the Wilson loop of the `Periodic Manifold' chapter (\cref{ch12_wilson_homotopy_lemma}).
\end{remark}

\begin{remark}[Strict boundary of causality]\label{rem:propagation_causality_boundary}
The causality given by \cref{cor:causality_speed} is a mathematical property of the \textbf{linear wave equation}; its derivation uses only the equation itself and energy conservation, and introduces no assumption beyond `linear + constant coefficients (or periodic coefficients) + finite energy'. In nonlinear media (strong-field optics, intense electromagnetic pulses in plasmas) the above conclusion of finite propagation speed \textbf{need not hold}---nonlinear terms can produce superluminal motion of the phase front. This chapter marks this boundary explicitly: all statements about propagation speed presuppose the \textbf{linear wave equation}.
\end{remark}

\section{Information evolution: Bayesian updating, activation functions and learning}\label{sec:information_evolution}

\begin{remark}[Positioning of this section and strict premises]\label{rem:information_evolution_positioning}
This section extends the mathematical framework of propagation and evolution from physical fields to \textbf{information states}. The core observation is that Bayesian updating, pointwise activation functions and parameter learning possess a unified algebraic structure within the translation-group framework already established in this chapter. All definitions, propositions, lemmas, theorems and corollaries in this section are bounded by the following two premises:

\begin{enumerate}[label=(\alph*)]
    \item \textbf{Linear evolution premise}: the evolution of an information state in \textbf{natural-parameter space} or \textbf{log-belief space} is generated by the action of the translation group $\mathbb{R}^m$ (structurally isomorphic to \cref{def:propagation_operator});
    \item \textbf{Pointwise activation premise}: the activation function is a \textbf{pointwise} map $\varphi:\mathbb{R}\to\mathbb{R}$ acting componentwise on the state vector and introducing no coupling between components.
\end{enumerate}

Any statement involving non-convex optimisation, adaptive learning rates, vector-valued activation, gating mechanisms and so on, which goes beyond the above premises, is invariably marked as a \textbf{note} or an \textbf{open problem}. The concepts of this section are the \textbf{mathematical foundation} of the corresponding concepts in the `Periodic Structure of Agents' chapter and the `Epicycle Activation Function and Function Approximation' chapter; the content of those two chapters becomes an application within the framework of this section.

Positioning of the seven themes of this section: translation---Bayesian updating in log-belief space is an additive translation (\cref{thm:bayesian_translation_structure}), and pointwise activation commutes with the propagation operator (\cref{prop:pointwise_activation_commutes}); scaling---the learning rate as the scaling factor of parametric evolution (\cref{def:learning_rate}); rotation---the phase rotation of the complex exponential activation function (\cref{prop:ceaf_phase_rotation}); conservation---monotone accumulation of the natural parameters under a conjugate prior (\cref{cor:natural_parameter_monotone}); symmetry---pointwise activation preserves translation invariance (\cref{cor:pointwise_preserves_translation}); closure---convergence closure of parametric evolution (\cref{thm:parametric_evolution_convergence}); period---the periodic structure of cosine-annealing learning rates (\cref{prop:cosine_annealing_period}).
\end{remark}

\subsection{The translational structure of Bayesian updating}\label{subsec:bayesian_translation}

\begin{definition}[Belief state and observation operator]\label{def:belief_state}
Let the parameter space be $\Theta\subseteq\mathbb{R}^m$. A \textbf{belief state} on $\Theta$ is a probability density $p:\Theta\to\mathbb{R}^+$ satisfying $\int_\Theta p(\theta)\,d\theta = 1$. The \textbf{observation operator} $\mathcal{O}_D$ maps a belief state $p$ and observed data $D$ to the posterior $p' = \mathcal{O}_D(p)$, defined by
\[
(\mathcal{O}_D p)(\theta) := \frac{p(D\mid\theta)\,p(\theta)}{Z_D}, \qquad Z_D := \int_\Theta p(D\mid\theta')\,p(\theta')\,d\theta'.
\]
If $Z_D = 0$, then $D$ and $p$ are said to be incompatible, and the observation operator is undefined there.
\end{definition}

\begin{definition}[Log-belief space]\label{def:log_belief_space}
Let $\mathcal{P}(\Theta)$ be the space of probability densities on $\Theta$. The \textbf{log-belief map} is defined by
\[
\Lambda:\mathcal{P}(\Theta)\to \mathcal{F}(\Theta), \qquad \Lambda(p) := \log p,
\]
where $\mathcal{F}(\Theta)$ is the space of real-valued functions on $\Theta$ (modulo constants). Two functions in $\mathcal{F}(\Theta)$ are regarded as identical if they differ by a constant. $\Lambda(p)$ is called the \textbf{log-belief} of $p$.
\end{definition}

\begin{theorem}[Translational structure of Bayesian updating]\label{thm:bayesian_translation_structure}
Let $p,p'\in\mathcal{P}(\Theta)$ with $p' = \mathcal{O}_D(p)$. Then in the log-belief space $\mathcal{F}(\Theta)$,
\[
\boxed{\ \Lambda(p') = \Lambda(p) + \ell_D,\ }
\]
where $\ell_D(\theta) := \log p(D\mid\theta) - \log Z_D$ is the \textbf{log-likelihood correction term} (modulo constants one may take $\log p(D\mid\theta)$). That is: \textbf{Bayesian updating in log-belief space is an additive translation with translation amount $\ell_D$}.
\end{theorem}

\begin{proof}
By \cref{def:belief_state} and \cref{def:log_belief_space},
\[
\Lambda(p')(\theta) = \log\Bigl(\frac{p(D\mid\theta)\,p(\theta)}{Z_D}\Bigr) = \log p(\theta) + \bigl[\log p(D\mid\theta) - \log Z_D\bigr] = \Lambda(p)(\theta) + \ell_D(\theta).
\]
In the sense of modulo constants, the $-\log Z_D$ part of $\ell_D$ may be omitted. This completes the proof.
\end{proof}

\begin{figure}[htbp]
\centering
\begin{tikzpicture}[>=Stealth, thick, node distance=1.3cm]
\node[draw, rounded corners, fill=blue!6, align=center] (p) {Belief state\\$p(\theta)$};
\node[draw, rounded corners, fill=orange!8, align=center, right=of p] (log) {Log-belief\\$\Lambda(p)=\log p$};
\node[draw, rounded corners, fill=green!8, align=center, right=of log] (post) {Posterior\\$p'=\mathcal{O}_D(p)$};
\node[draw, rounded corners, fill=purple!8, align=center, below=of log] (shift) {Additive translation\\$\Lambda(p')=\Lambda(p)+\ell_D$};
\draw[->] (p) -- node[above] {$\Lambda$} (log);
\draw[->] (p) -- node[below] {Bayesian updating} (post);
\draw[->] (log) -- (shift);
\draw[->] (shift) -- (post);
\node at (5.5,-2.0) {Translation amount $\ell_D(\theta)=\log p(D\mid\theta)-\log Z_D$.};
\end{tikzpicture}
\caption{Additive translation in log-belief space.}
\label{fig:bayesian_translation}
\end{figure}

\begin{corollary}[Structural isomorphism between Bayesian updating and the propagation operator]\label{cor:bayesian_propagation_isomorphism}
Bayesian updating (\cref{thm:bayesian_translation_structure}) and the propagation operator (\cref{def:propagation_operator}) are structurally isomorphic:
\[
\begin{array}{c|c|c}
& \text{action space} & \text{form of action} \\ \hline
\text{propagation operator} & L^2(\mathbb{R}^n) & f(\mathbf{x})\mapsto f(\mathbf{x}-\mathbf{a}) \\
\text{Bayesian updating} & \mathcal{F}(\Theta) & \Lambda(p)(\theta)\mapsto \Lambda(p)(\theta)+\ell_D(\theta)
\end{array}
\]
Their common algebraic structure is a \textbf{translation-group action}: the propagation operator is generated by the spatial translation group $\mathbb{R}^n$, whilst Bayesian updating is generated by the \textbf{data-dependent translation} determined by the log-likelihood.
\end{corollary}

\begin{proof}
Compare \cref{prop:propagation_group}(i) (the propagation operators form a group) directly with \cref{thm:bayesian_translation_structure} (Bayesian updating is an additive translation).
\end{proof}

\begin{definition}[Exponential family and natural parameters]\label{def:exponential_family}
Let $\mathcal{X}$ be the data space and $\Theta\subseteq\mathbb{R}^m$ the parameter space. A likelihood family $\{p(\cdot\mid\theta)\}_{\theta\in\Theta}$ is called an \textbf{exponential family} if there exist a sufficient statistic $T:\mathcal{X}\to\mathbb{R}^d$, a natural-parameter map $\eta:\Theta\to\mathbb{R}^d$, a log-partition function $A:\Theta\to\mathbb{R}$ and a base measure $h:\mathcal{X}\to\mathbb{R}^+$ such that
\[
p(x\mid\theta) = h(x)\exp\bigl(\eta(\theta)\cdot T(x) - A(\theta)\bigr), \qquad x\in\mathcal{X}.
\]
$\eta(\Theta)\subseteq\mathbb{R}^d$ is called the \textbf{natural-parameter space}.
\end{definition}

\begin{definition}[Conjugate prior and natural-parameter coordinates]\label{def:conjugate_prior_natural}
Let $\{p(\cdot\mid\theta)\}$ be an exponential family (\cref{def:exponential_family}). The \textbf{conjugate prior family} is
\[
\pi(\theta\mid\tau,\nu) := \frac{\exp\bigl(\tau\cdot\eta(\theta) - \nu A(\theta)\bigr)}{\int_\Theta \exp\bigl(\tau\cdot\eta(\theta') - \nu A(\theta')\bigr)\,d\theta'}, \qquad \tau\in\mathbb{R}^d,\ \nu>0.
\]
$(\tau,\nu)$ are called the \textbf{natural-parameter coordinates} of the conjugate prior. When $d=1$, $(\tau,\nu)\in\mathbb{R}\times\mathbb{R}^+$.
\end{definition}

\begin{theorem}[Exact finite-order evolution of Bayesian updating under a conjugate prior]\label{thm:conjugate_exact_finite_order}
Let $\{p(\cdot\mid\theta)\}$ be an exponential family (\cref{def:exponential_family}), let the prior be the conjugate prior $\pi(\theta\mid\tau_0,\nu_0)$ (\cref{def:conjugate_prior_natural}), and let the observation sequence $x_1,\dots,x_n$ be independent and identically distributed from $p(\cdot\mid\theta)$. Then the posterior is again a conjugate prior, with natural-parameter coordinates
\[
\boxed{\ \tau_n = \tau_0 + \sum_{i=1}^n T(x_i), \qquad \nu_n = \nu_0 + n.\ }
\]
That is: \textbf{Bayesian updating under a conjugate prior is, in the natural-parameter coordinates $(\tau,\nu)$, a finite-order additive evolution along the fixed direction $(T(x_i),1)$}---each observation corresponds to one translation.
\end{theorem}

\begin{proof}
By Bayes' formula and the definition of the conjugate prior:
\[
\pi(\theta\mid x_1,\dots,x_n) \propto \pi(\theta\mid\tau_0,\nu_0)\prod_{i=1}^n p(x_i\mid\theta) \propto \exp\Bigl(\bigl(\tau_0 + \textstyle\sum_i T(x_i)\bigr)\cdot\eta(\theta) - (\nu_0+n)A(\theta)\Bigr).
\]
Comparing with the conjugate prior form (\cref{def:conjugate_prior_natural}) gives the result. Finite order: in the $(\tau,\nu)$ coordinates each update is $(\tau,\nu)\mapsto(\tau+T(x_i), \nu+1)$---a first-order affine transformation.
\end{proof}

\begin{corollary}[Monotone accumulation of natural parameters]\label{cor:natural_parameter_monotone}
In the notation of \cref{thm:conjugate_exact_finite_order}, $\nu_n$ is strictly increasing in $n$, and $\tau_n$ accumulates along the direction $T(x_i)$. This monotonicity is the algebraic root of the posterior contraction in the Bernstein--von Mises theorem (\cref{thm:bernstein_von_mises}) of the `Periodic Structures in Probabilistic Systems' chapter.
\end{corollary}

\begin{proof}
$\nu_n = \nu_0+n$ is strictly increasing (it increases by exactly $1$ at each step). $\tau_n - \tau_{n-1} = T(x_n)$ is the sufficient statistic of the $n$-th observation.
\end{proof}

\begin{example}[Exact evolution of Beta--binomial conjugacy]\label{ex:beta_binomial_exact}
Let the likelihood be $X\mid\theta\sim\operatorname{Bernoulli}(\theta)$, with exponential-family form $p(x\mid\theta) = \theta^x(1-\theta)^{1-x}$, $T(x)=x$, $\eta(\theta) = \log(\theta/(1-\theta))$ and $A(\theta) = -\log(1-\theta)$. The conjugate prior is the Beta distribution $\operatorname{Beta}(\alpha,\beta)$, with natural-parameter coordinates $(\tau,\nu) = (\alpha-1, \alpha+\beta-2)$ (the precise choice of coordinates depends on convention). After $n$ observations,
\[
\alpha_n = \alpha_0 + \sum_{i=1}^n x_i, \qquad \beta_n = \beta_0 + n - \sum_{i=1}^n x_i.
\]
This agrees with the conclusion of the Beta--binomial conjugate model in the `Periodic Structures in Probabilistic Systems' chapter (\cref{ex:beta_binomial}); here its \textbf{translation structure} is given in natural-parameter coordinates.
\end{example}

\begin{example}[Exact evolution of Gaussian--Gaussian conjugacy]\label{ex:gaussian_gaussian_exact}
Let the likelihood be $X\mid\mu\sim\mathcal{N}(\mu,\sigma^2)$ (with $\sigma^2$ known) and let the prior on the parameter $\mu$ be $\mathcal{N}(\mu_0,\tau_0^2)$. Taking the precision-mean product coordinates $(1/\tau^2,\ \mu/\tau^2)$ as natural-parameter coordinates, after $n$ observations
\[
\frac{1}{\tau_n^2} = \frac{1}{\tau_0^2} + \frac{n}{\sigma^2}, \qquad \frac{\mu_n}{\tau_n^2} = \frac{\mu_0}{\tau_0^2} + \frac{\sum_i x_i}{\sigma^2}.
\]
This agrees with the normal--normal conjugate model of the `Periodic Structures in Probabilistic Systems' chapter (\cref{ex:normal_normal}); here its \textbf{translation structure} is given in natural-parameter coordinates.
\end{example}

\begin{remark}[Correspondence table of finite-order evolution and the propagation operator]\label{rem:belief_propagation_correspondence}
Combining \cref{thm:bayesian_translation_structure} with \cref{thm:conjugate_exact_finite_order}, Bayesian updating and the propagation operator satisfy the following structural correspondence:
\begin{center}
\renewcommand{\arraystretch}{1.3}
\begin{tabular}{@{}p{3.2cm} p{4.2cm} p{4.2cm}@{}}
\toprule
\textbf{Object} & \textbf{Propagation operator} & \textbf{Bayesian updating} \\
\midrule
Action space & field function $L^2(\mathbb{R}^n)$ & log-belief $\mathcal{F}(\Theta)$ or natural-parameter space $\mathbb{R}^d\times\mathbb{R}^+$ \\
Translation amount & spatial vector $\mathbf{a}$ & log-likelihood $\ell_D$ or sufficient statistic $T(x)$ \\
Generator & $-\nabla$ (spatial directional derivative) & observation operator (integration in the likelihood direction) \\
Group structure & translation group $\mathbb{R}^n$ & data-dependent translation (each observation defines one translation) \\
Transfer matrix (frequency domain) & $e^{-i\mathbf{k}\cdot\mathbf{a}}$ & addition of the natural parameters under a conjugate prior \\
\bottomrule
\end{tabular}
\end{center}
\end{remark}

\subsection{Pointwise activation functions and the propagation operator}\label{subsec:pointwise_activation}

\begin{definition}[Pointwise activation function]\label{def:pointwise_activation}
A map $\varphi:\mathbb{R}\to\mathbb{R}$ is called a \textbf{pointwise activation function}. Its action on a function $f:\mathbb{R}^n\to\mathbb{R}$ is defined by
\[
(\varphi\cdot f)(\mathbf{x}) := \varphi(f(\mathbf{x})), \qquad \mathbf{x}\in\mathbb{R}^n.
\]
The family of pointwise activation functions includes: the identity map $\varphi(x)=x$, the Heaviside step $H(x)$, ReLU $\max(0,x)$, Sigmoid $\sigma(x)$, Tanh, GELU, Swish, and complex exponential activation (\cref{def:ceaf_definition}).
\end{definition}

\begin{proposition}[Commutativity of pointwise activation and the propagation operator]\label{prop:pointwise_activation_commutes}
Let $\varphi:\mathbb{R}\to\mathbb{R}$ be a pointwise activation function (\cref{def:pointwise_activation}) and let $P_{\mathbf{a}}$ be a propagation operator (\cref{def:propagation_operator}). Then
\[
\boxed{\ P_{\mathbf{a}}\circ\varphi = \varphi\circ P_{\mathbf{a}}.\ }
\]
That is, a pointwise activation function and the propagation operator commute exactly.
\end{proposition}

\begin{proof}
For arbitrary $f:\mathbb{R}^n\to\mathbb{R}$ and $\mathbf{x}\in\mathbb{R}^n$, compute both sides:
\[
(P_{\mathbf{a}}\circ\varphi)(f)(\mathbf{x}) = (\varphi\cdot f)(\mathbf{x}-\mathbf{a}) = \varphi(f(\mathbf{x}-\mathbf{a})),
\]
\[
(\varphi\circ P_{\mathbf{a}})(f)(\mathbf{x}) = \varphi((P_{\mathbf{a}}f)(\mathbf{x})) = \varphi(f(\mathbf{x}-\mathbf{a})).
\]
The two sides are equal.
\end{proof}

\begin{figure}[htbp]
\centering
\begin{tikzpicture}[>=Stealth, thick, node distance=2.2cm]
\node (f) {$f$};
\node[right=of f] (pf) {$P_{\mathbf{a}}f$};
\node[below=of f] (phif) {$\varphi\cdot f$};
\node[below=of pf] (papphi) {$P_{\mathbf{a}}(\varphi\cdot f)$};
\draw[->] (f) -- node[above] {$P_{\mathbf{a}}$} (pf);
\draw[->] (f) -- node[left] {$\varphi$} (phif);
\draw[->] (pf) -- node[right] {$\varphi$} (papphi);
\draw[->] (phif) -- node[below] {$P_{\mathbf{a}}$} (papphi);
\node at (1.2,-2.4) {$P_{\mathbf{a}}\circ\varphi=\varphi\circ P_{\mathbf{a}}$: pointwise activation and the propagation operator commute exactly.};
\end{tikzpicture}
\caption{Commutative diagram for pointwise activation and the propagation operator.}
\label{fig:pointwise_commute}
\end{figure}

\begin{corollary}[Pointwise activation preserves translation invariance]\label{cor:pointwise_preserves_translation}
Let $f$ be a $P_{\mathbf{a}}$-invariant function (that is, $P_{\mathbf{a}}f = f$). Then $\varphi\cdot f$ is also $P_{\mathbf{a}}$-invariant:
\[
P_{\mathbf{a}}(\varphi\cdot f) = \varphi\cdot f.
\]
\end{corollary}

\begin{proof}
By \cref{prop:pointwise_activation_commutes}, $P_{\mathbf{a}}(\varphi\cdot f) = \varphi\cdot(P_{\mathbf{a}}f) = \varphi\cdot f$.
\end{proof}

\begin{remark}[Linear transformations break translation invariance]\label{rem:linear_transform_breaks_translation}
Pointwise activation commutes with the propagation operator (\cref{prop:pointwise_activation_commutes}), but a general \textbf{linear transformation} $W:\mathbb{R}^n\to\mathbb{R}^n$ does not commute with the propagation operator: $W\circ P_{\mathbf{a}}\neq P_{\mathbf{a}}\circ W$ (unless $W$ is compatible with the direction of $\mathbf{a}$). This explains the algebraic necessity of the alternating structure `linear transformation $\to$ pointwise activation $\to$ linear transformation $\to\cdots$' in neural networks: a linear transformation breaks translation invariance, a pointwise activation preserves it, and the alternation of the two produces intermediate representations that are neither completely invariant nor completely arbitrary.
\end{remark}

\begin{definition}[Complex exponential activation function]\label{def:ceaf_definition}
Let $N\in\mathbb{N}^+$, $A_k>0$, $\omega_k\in\mathbb{R}$ and $\phi_k\in[0,2\pi)$. The \textbf{complex exponential activation function} (CEAF) is defined by
\[
\mathcal{E}(x) := \sum_{k=1}^{N}A_k e^{i(\omega_k x+\phi_k)}, \qquad x\in\mathbb{R}.
\]
Its readout modes include: the real-part readout $\operatorname{Re}\mathcal{E}(x)$, the two-channel readout $(\operatorname{Re}\mathcal{E}(x),\operatorname{Im}\mathcal{E}(x))$, and the modulus-phase readout $(|\mathcal{E}(x)|, \arg\mathcal{E}(x))$.
\end{definition}

\begin{proposition}[Propagation structure of the CEAF]\label{prop:ceaf_phase_rotation}
Under translation the CEAF behaves as a pure phase rotation: for every $a\in\mathbb{R}$,
\[
\mathcal{E}(x+a) = \sum_{k=1}^{N}A_k e^{i(\omega_k x+\phi_k+\omega_k a)}.
\]
That is, a translation by $a$ is equivalent to increasing the phase of the $k$-th component by $\omega_k a$---this is the componentwise realisation on the CEAF of \cref{prop:phase_propagation} (consistency of the propagation operator and the phase operator).
\end{proposition}

\begin{proof}
Substitute \cref{def:ceaf_definition} and $e^{i(\omega_k(x+a)+\phi_k)} = e^{i(\omega_k x+\phi_k)}e^{i\omega_k a}$.
\end{proof}

\begin{corollary}[Periodicity condition of the CEAF]\label{cor:ceaf_period}
If $\omega_k/\omega_0\in\mathbb{Z}$ for all $k$ (where $\omega_0$ is the base frequency), then $\mathcal{E}$ is periodic with period $T = 2\pi/\omega_0$; otherwise $\mathcal{E}$ is an almost periodic function. This is the term-by-term realisation on the CEAF of the periodic-function criterion of the `Operations and Periodicity' chapter.
\end{corollary}

\begin{proof}
If $\omega_k = n_k\omega_0$ ($n_k\in\mathbb{Z}$), then $\mathcal{E}(x+T) = \sum_k A_k e^{i(n_k\omega_0 x+\phi_k+n_k\omega_0\cdot 2\pi/\omega_0)} = \sum_k A_k e^{i(n_k\omega_0 x+\phi_k)} = \mathcal{E}(x)$. If there exists an incommensurable pair of frequencies, then $\mathcal{E}$ is a Bohr almost periodic function.
\end{proof}

\begin{example}[Phase rotation of a single-component CEAF]\label{ex:single_ceaf_phase}
Take $\mathcal{E}(x) = e^{i\omega x}$ ($N=1$, $A_1=1$, $\phi_1=0$). A translation by $a$ gives $\mathcal{E}(x+a) = e^{i\omega(x+a)} = e^{i\omega x}e^{i\omega a}$---the phase increases by $\omega a$. When $a = 2\pi/\omega$ the phase increases by $2\pi$ and $\mathcal{E}$ returns to its starting value---this is the periodic closure of propagation on the CEAF (\cref{cor:propagation_period_closure}).
\end{example}

\begin{remark}[Relation between the CEAF and kernel-smoothed derivative-integral descent]\label{rem:ceaf_kdi_interface}
The parameters $\{A_k,\omega_k,\phi_k\}$ of the CEAF are optimised by kernel-smoothed derivative-integral descent (\cref{def:kdi_definition}). This section establishes only the propagation and phase structure of the CEAF (\cref{prop:ceaf_phase_rotation}, \cref{cor:ceaf_period}); its optimisation algorithm is given in \cref{subsec:parametric_evolution} as a special case of parametric evolution.
\end{remark}

\subsection{Learning as parametric evolution}\label{subsec:parametric_evolution}

\begin{definition}[Parametric evolution equation]\label{def:parametric_evolution}
Let the loss function $\mathcal{L}:\mathbb{R}^d\to\mathbb{R}$ be $C^1$ smooth and let $\eta>0$ be the \textbf{learning rate} (scaling factor). The \textbf{parametric evolution equation} is
\[
\theta_{t+1} = \theta_t - \eta\,\nabla\mathcal{L}(\theta_t), \qquad t = 0,1,2,\dots.
\]
If $\mathcal{L}$ attains a minimum at $\theta^*$, then the behaviour of $\theta_t$ as $t\to\infty$ is called \textbf{parametric evolution}.
\end{definition}

\begin{definition}[Learning rate]\label{def:learning_rate}
The quantity $\eta>0$ in \cref{def:parametric_evolution} is called the \textbf{learning rate}. The learning rate is the scaling factor of parametric evolution---the larger $\eta$ is, the larger the parameter translation at each step (fast but possibly unstable); the smaller $\eta$ is, the smaller the parameter translation at each step (slow but more stable).
\end{definition}

\begin{theorem}[Linear evolution and convergence condition under quadratic loss]\label{thm:parametric_evolution_convergence}
Let the loss function be the quadratic form $\mathcal{L}(\theta) = \frac{1}{2}(\theta-\theta^*)^{\mathsf T}H(\theta-\theta^*)$, where $H$ is a symmetric positive definite matrix with eigenvalues $0 < \lambda_1\le\dots\le\lambda_d$. Then parametric evolution (\cref{def:parametric_evolution}) is the \textbf{linear evolution}
\[
\theta_{t+1} - \theta^* = (I - \eta H)(\theta_t - \theta^*),
\]
and the error along the $i$-th eigendirection is
\[
e_t^{(i)} = (1-\eta\lambda_i)^t\,e_0^{(i)}.
\]
The \textbf{necessary and sufficient condition for convergence} is
\[
0 < \eta < \frac{2}{\lambda_d}.
\]
When $\eta = 2/(\lambda_1+\lambda_d)$, the maximum of the convergence rates $|1-\eta\lambda_i|$ is least---in that case $\max_i|1-\eta\lambda_i| = (\lambda_d-\lambda_1)/(\lambda_1+\lambda_d)$.
\end{theorem}

\begin{proof}
Decompose the error vector $\mathbf{e}_t := \theta_t - \theta^* = \sum_i e_t^{(i)}\mathbf{v}_i$ along the eigenbasis of $H$. Since $\nabla\mathcal{L}(\theta_t) = H(\theta_t-\theta^*)$, the parametric evolution equation gives $\theta_{t+1}-\theta^* = (I-\eta H)(\theta_t-\theta^*)$. The eigendirections decouple: $e_{t+1}^{(i)} = (1-\eta\lambda_i)e_t^{(i)}$, and iteration gives $e_t^{(i)} = (1-\eta\lambda_i)^t e_0^{(i)}$. Convergence holds if and only if $|1-\eta\lambda_i|<1$ for every $i$, that is, $0<\eta\lambda_i<2$, which is equivalent to $0<\eta<2/\lambda_d$. The optimal learning rate is given by minimising the spectral radius $\rho(\eta) = \max_i|1-\eta\lambda_i|$; the spectral radius is least at $\eta = 2/(\lambda_1+\lambda_d)$, where the convergence rates in the $\lambda_1$ and $\lambda_d$ directions are equal, namely $(\lambda_d-\lambda_1)/(\lambda_1+\lambda_d)$.
\end{proof}

\begin{figure}[htbp]
\centering
\begin{tikzpicture}[>=Stealth, thick, scale=0.9]
\draw[->] (-0.5,0) -- (5.5,0) node[right] {$\theta_1$};
\draw[->] (0,-0.5) -- (0,4.5) node[above] {$\theta_2$};
\draw[rotate=35, blue, very thick] (0,0) ellipse (2.8 and 1.2);
\draw[rotate=35, red, thick, dashed] (0,0) ellipse (2.0 and 0.85);
\draw[rotate=35, green!60!black, thick, dotted] (0,0) ellipse (1.2 and 0.5);
\fill (0,0) circle (2pt) node[below right] {$\theta^*$};
\draw[->, red, very thick] (4.3,3.6) .. controls (3.0,2.5) and (1.8,1.6) .. (0.2,0.2);
\node[red] at (4.5,3.9) {Trajectory of $\theta_t$};
\node at (2.8,-0.9) {Along an eigendirection: $e_t^{(i)}=(1-\eta\lambda_i)^t e_0^{(i)}$};
\node at (2.8,-1.6) {Convergence condition: $0<\eta<2/\lambda_d$};
\end{tikzpicture}
\caption{Convergence of parametric evolution under quadratic loss.}
\label{fig:parametric_convergence}
\end{figure}

\begin{corollary}[Parametric evolution as the discrete realisation of the linear evolution semigroup]\label{cor:parametric_as_semigroup}
Parametric evolution (\cref{thm:parametric_evolution_convergence}) is the discrete realisation of the evolution semigroup of this chapter (\cref{def:evolution_semigroup}): the generator is $-H$ and the semigroup is $\{(I-\eta H)^t\}_{t\ge 0}$. In the continuum limit $\eta\to 0$,
\[
\theta(t) = e^{-Ht}\theta(0)
\]
corresponds to continuous-time parametric evolution.
\end{corollary}

\begin{proof}
Iterating the parametric evolution equation $\theta_{t+1} = (I-\eta H)\theta_t$ gives $\theta_t = (I-\eta H)^t\theta_0$. Taking $\eta\to 0$, $t\to\infty$ and $\eta t\to s$, we get $(I-\eta H)^t\to e^{-Hs}$, which is the continuum limit.
\end{proof}

\begin{definition}[Kernel-smoothed derivative-integral descent]\label{def:kdi_definition}
Let the loss function $\mathcal{L}:\mathbb{R}^d\to\mathbb{R}$ be $C^1$ smooth and let $\theta\in\mathbb{R}^d$ be a parameter. The update rule of \textbf{kernel-smoothed derivative-integral descent} (KDI) is
\[
\theta_{t+1} = \theta_t - \eta\Bigl[\alpha\,\nabla\mathcal{L}(\theta_t) + \beta\,\mathbf{d}_{\text{int}}(\theta_t) + \gamma\,\mathbf{d}_{\text{curv}}(\theta_t)\Bigr],
\]
where $\alpha,\beta,\gamma\ge 0$ are weight coefficients, $\mathbf{d}_{\text{int}}$ is the \textbf{accumulated-variation correction term} (integral information) and $\mathbf{d}_{\text{curv}}$ is the \textbf{curvature correction term} (second-derivative information). When $\beta=\gamma=0$, KDI degenerates to standard gradient descent.
\end{definition}

\begin{proposition}[Descent property of the KDI direction]\label{prop:kdi_descent}
Let $\mathcal{L}$ be $C^1$ in a neighbourhood of $\theta$ and let the KDI direction be
\[
\mathbf{d}_{\text{KDI}}(\theta) := \alpha\,\nabla\mathcal{L}(\theta) + \beta\,\mathbf{d}_{\text{int}}(\theta) + \gamma\,\mathbf{d}_{\text{curv}}(\theta).
\]
If $\nabla\mathcal{L}(\theta)\cdot\mathbf{d}_{\text{KDI}}(\theta) > 0$ (that is, the inner product of the KDI direction with the negative-gradient direction is positive), then there exists $\eta>0$ such that
\[
\mathcal{L}(\theta - \eta\,\mathbf{d}_{\text{KDI}}(\theta)) < \mathcal{L}(\theta).
\]
\end{proposition}

\begin{proof}
The directional derivative is $\frac{d}{d\eta}\big|_{\eta=0}\mathcal{L}(\theta-\eta\mathbf{d}_{\text{KDI}}) = -\nabla\mathcal{L}(\theta)\cdot\mathbf{d}_{\text{KDI}}(\theta)$. If this value is negative (that is, $\nabla\mathcal{L}\cdot\mathbf{d}_{\text{KDI}} > 0$), then the function value decreases strictly for sufficiently small $\eta$.
\end{proof}

\begin{remark}[Strict boundary of KDI]\label{rem:kdi_boundary}
\cref{prop:kdi_descent} states only the \textbf{descent property} of KDI---as long as the angle between the direction and the negative gradient is less than $90^\circ$, descent occurs. The \textbf{advantages} of KDI over standard gradient descent (escaping saddle points, crossing local minima, accelerating across flat regions) depend on the geometric structure of the specific problem and \textbf{do not constitute theorems of this chapter}. The convergence analysis of KDI (especially in the non-convex case) is an \textbf{open problem} (\cref{rem:information_evolution_open}).
\end{remark}

\begin{definition}[Learning-rate schedule]\label{def:learning_rate_schedule}
A \textbf{learning-rate schedule} is the learning rate $\eta_t$ regarded as a function of the iteration step $t$. Common types:
\begin{enumerate}[label=(\roman*)]
    \item \textbf{Step decay}: $\eta_t = \eta_0\gamma^{\lfloor t/T\rfloor}$ ($\gamma\in(0,1)$, $T$ is the decay period);
    \item \textbf{Cosine annealing}: $\eta_t = \eta_{\min} + \frac{1}{2}(\eta_{\max}-\eta_{\min})\bigl(1+\cos(\pi t/T)\bigr)$;
    \item \textbf{Exponential decay}: $\eta_t = \eta_0 e^{-t/\tau}$.
\end{enumerate}
A learning-rate schedule is a modulation in time of the scaling factor of parametric evolution---from a large learning rate (rapid exploration) to a small learning rate (fine convergence).
\end{definition}

\begin{proposition}[Periodic structure of the cosine-annealing learning rate]\label{prop:cosine_annealing_period}
The cosine-annealing learning rate $\eta_t = \eta_{\min} + \frac{1}{2}(\eta_{\max}-\eta_{\min})(1+\cos(\pi t/T))$ is a periodic function of period $2T$:
\[
\eta_{t+2T} = \eta_t, \qquad \forall t\ge 0.
\]
Within each period, $\eta_t$ decreases monotonically from $\eta_{\max}$ to $\eta_{\min}$ ($t\in[0,T]$), and then increases monotonically back to $\eta_{\max}$ ($t\in[T,2T]$).
\end{proposition}

\begin{proof}
Periodicity follows from $\cos(\pi(t+2T)/T) = \cos(\pi t/T + 2\pi) = \cos(\pi t/T)$. The monotonicity is given by the monotonicity of the cosine on $[0,\pi]$ and $[\pi,2\pi]$.
\end{proof}

\begin{figure}[htbp]
\centering
\begin{tikzpicture}[>=Stealth, thick, scale=0.9]
\draw[->] (-0.3,0) -- (9,0) node[right] {$t$};
\draw[->] (0,-0.2) -- (0,3.2) node[above] {$\eta_t$};
\draw[blue, very thick, domain=0:8, samples=200]
  plot (\x, {0.4 + 1.2*(1+cos(180*\x/4))/2});
\draw[dashed, gray] (0,1.6) -- (8,1.6) node[right] {$\eta_{\max}$};
\draw[dashed, gray] (0,0.4) -- (8,0.4) node[right] {$\eta_{\min}$};
\draw[<->] (4,0.1) -- (8,0.1) node[midway, below] {Period $2T$};
\node at (4.5,2.6) {$\eta_t=\eta_{\min}+\frac12(\eta_{\max}-\eta_{\min})(1+\cos(\pi t/T))$};
\end{tikzpicture}
\caption{Periodic structure of the cosine-annealing learning rate.}
\label{fig:cosine_annealing}
\end{figure}

\begin{example}[Numerical verification of parametric evolution]\label{ex:parametric_evolution_numeric}
Take $d=2$, $H = \operatorname{diag}(1,4)$, $\theta^* = (0,0)$ and the initial value $\theta_0 = (1,1)$. By \cref{thm:parametric_evolution_convergence}, the convergence condition is $0 < \eta < 2/4 = 0.5$. Take $\eta = 0.3$:
\begin{itemize}
    \item First eigendirection: $|1-\eta\lambda_1| = |1-0.3| = 0.7$;
    \item Second eigendirection: $|1-\eta\lambda_2| = |1-1.2| = 0.2$.
\end{itemize}
The error $\|\theta_t - \theta^*\|$ converges with a decay factor of at least $0.7$ per step. Take $\eta = 0.6 > 0.5$: $|1-\eta\lambda_2| = |1-2.4| = 1.4 > 1$, and the second eigendirection diverges---parametric evolution is unstable.
\end{example}

\begin{theorem}[Convergence closure of parametric evolution]\label{thm:parametric_evolution_closure}
Let $\mathcal{L}$ be an $L$-smooth, $\mu$-strongly convex function on $\mathbb{R}^d$ ($0<\mu\le L$). Then with $\eta = 1/L$, parametric evolution (\cref{def:parametric_evolution}) satisfies
\[
\|\theta_t - \theta^*\|^2 \le \Bigl(1-\frac{\mu}{L}\Bigr)^t\,\|\theta_0-\theta^*\|^2.
\]
In particular, $\theta_t\to\theta^*$ converges exponentially---this is the \textbf{convergence closure of parametric evolution}.
\end{theorem}

\begin{proof}
By $L$-smoothness (the descent lemma) and $\mu$-strong convexity, the standard analysis gives
\[
\mathcal{L}(\theta_{t+1}) - \mathcal{L}(\theta^*) \le \Bigl(1-\frac{\mu}{L}\Bigr)\bigl(\mathcal{L}(\theta_t)-\mathcal{L}(\theta^*)\bigr).
\]
By $\mu$-strong convexity, $\|\theta_t-\theta^*\|^2\le\frac{2}{\mu}(\mathcal{L}(\theta_t)-\mathcal{L}(\theta^*))$; substituting gives the result.
\end{proof}

\begin{example}[Sequential updating of a conjugate prior and composition of the propagation operator]\label{ex:sequential_bayesian_propagation}
Consider the Beta--binomial conjugate model (\cref{ex:beta_binomial_exact}). Each observation $x_i\in\{0,1\}$ corresponds to one translation of the natural parameters:
\[
(\tau,\nu)\mapsto(\tau + T(x_i),\nu+1) = (\tau+x_i,\nu+1).
\]
The accumulation over $n$ observations is \cref{thm:conjugate_exact_finite_order}. This is a concrete instance of \textbf{Bayesian updating as a composition of propagation operators}---each observation corresponds to a translation operator $P_{(\tau+x_i,\nu+1)}$, and $n$ observations correspond to the composition $P_{(\tau+\sum x_i,\nu+n)}$.
\end{example}

\begin{remark}[Unified table of the three information evolutions]\label{rem:three_info_evolutions}
The three information evolutions established in this section (Bayesian updating, pointwise activation and parameter learning) are unified within the framework of this chapter as translation-group actions:
\begin{center}
\small
\renewcommand{\arraystretch}{1.3}
\setlength{\tabcolsep}{4pt}
\begin{tabular}{@{}p{2.2cm} p{3.6cm} p{3.6cm} p{2.2cm}@{}}
\toprule
\textbf{Information evolution} & \textbf{Action space} & \textbf{Translation amount} & \textbf{Group generator} \\
\midrule
Bayesian updating & log-belief space $\mathcal{F}(\Theta)$ & log-likelihood $\ell_D$ & observation operator \\
Conjugate-prior Bayesian updating & natural-parameter space $\mathbb{R}^d\times\mathbb{R}^+$ & sufficient statistic $(T(x),1)$ & data-dependent translation \\
Pointwise activation & function space (pointwise action) & no translation (commutes with $P_{\mathbf{a}}$) & pointwise map \\
Parameter learning (quadratic loss) & parameter space $\mathbb{R}^d$ & translation in the gradient direction & $-H$ \\
\bottomrule
\end{tabular}
\end{center}
\end{remark}

\begin{remark}[Interfaces with other chapters of the book]\label{rem:information_evolution_interfaces}
The interfaces between the mathematical structures established in this section and the other chapters of the book are as follows. All the entries cited from the `Periodic Structure of Agents' chapter and the `Epicycle Activation Function and Function Approximation' chapter \textbf{come after this chapter}---the definitions and propositions of this section are the \textbf{mathematical foundation} of the corresponding concepts in those chapters, and the content of those chapters becomes an application within the framework of this section.

\begin{enumerate}[label=(\roman*)]
    \item \textbf{The `Periodic Structures in Probabilistic Systems' chapter}: \cref{thm:bayesian_translation_structure} and \cref{thm:conjugate_exact_finite_order} provide the algebraic foundation of the translation-group structure for sequential Bayesian updating (\cref{def:sequential_bayesian}) and conjugate priors (\cref{def:conjugate_prior});
    \item \textbf{The `Periodic Structure of Agents' chapter}: the CEAF (\cref{def:ceaf_definition}) and parametric evolution (\cref{def:parametric_evolution}) provide the mathematical foundation for that chapter's neural-network training dynamics; \cref{def:eaf_agent} and \cref{def:kdi_agent} in that chapter become, within the framework of this section, application instances of \cref{def:ceaf_definition} and \cref{def:kdi_definition};
    \item \textbf{The `Epicycle Activation Function and Function Approximation' chapter}: the propagation structure of the CEAF (\cref{prop:ceaf_phase_rotation}) provides a phase-analysis tool for that chapter's approximation theory;
    \item \textbf{The `Groups and Group Functions' chapter}: \cref{prop:pointwise_activation_commutes} is the direct application to activation functions of that chapter's commutation law between pointwise operations and group actions (\cref{ch8_pointwise_invariance}).
\end{enumerate}
\end{remark}

\begin{remark}[Open problems of this section]\label{rem:information_evolution_open}
The following problems lie beyond the range of the theorems already established in this section and are listed as open problems \textbf{of this section}:
\begin{enumerate}[label=(\roman*)]
    \item \textbf{Finite-order evolution under a non-conjugate prior}: when the prior is not a conjugate prior, Bayesian updating generally \textbf{cannot} be represented as a finite-order translation in a finite-dimensional natural-parameter space. Whether there exists a more general finite-dimensional embedding that makes the update a finite-order evolution is an open problem;
    \item \textbf{Convergence of parametric evolution under a non-convex loss}: both \cref{thm:parametric_evolution_convergence} and \cref{thm:parametric_evolution_closure} require $\mathcal{L}$ to be convex. Under a non-convex loss, parametric evolution may become trapped in a local minimum or a saddle point, and its convergence behaviour (in particular the global convergence of KDI) is an open problem;
    \item \textbf{Commutation of propagation under vector-valued activation}: \cref{prop:pointwise_activation_commutes} relies on the pointwise premise. The commutation of a vector-valued activation (such as softmax) with the propagation operator requires additional conditions and lies beyond the scope of this section;
    \item \textbf{Group structure under adaptive learning rates}: in \cref{def:learning_rate_schedule} the learning-rate schedule depends on the iteration step $t$. For an adaptive learning rate (such as Adam) the learning rate depends on the gradient history, and parametric evolution is then no longer a representation of a fixed group---its algebraic structure is an open problem.
\end{enumerate}
\end{remark}

\begin{remark}[Open problems of the whole chapter]\label{rem:propagation_open}
The following problems lie beyond the range of the theorems already established in this chapter and are listed as open problems \textbf{of the chapter as a whole} (to be distinguished from the section-level open problems of \cref{rem:information_evolution_open}---the latter belong to Section \ref{sec:information_evolution}, whilst the former belong to the remaining sections of this chapter taken as a whole):
\begin{enumerate}[label=(\roman*)]
    \item a rigorous dispersion theory of wave packets in nonlinear periodic media;
    \item Anderson localisation of electromagnetic wave packets in random periodic media, and a quantitative description of their propagation;
    \item a rigorous propagation theory of the Airy kernel near a caustic surface (this chapter only lists it in \cref{rem:caustic_eikonal_failure} and in the kernel table \ref{tab:all_kernels}, without developing it);
    \item the extension of the propagation operator and of periodic convolution to non-commutative groups (such as $SU(2)$ or more general Lie groups), and its interface with the phase functor of the `Periodic Category' chapter;
    \item the doubly periodic structure of propagation together with time translation and spatial translation in time-varying periodic media.
\end{enumerate}
The solution of these problems would advance the theory of propagation and evolution from the linear constant-coefficient case to more general periodic systems.
\end{remark}

\begin{remark}[Seven-theme correspondence of this chapter with the whole book]\label{rem:propagation_seven_themes}
\begin{center}
\small
\renewcommand{\arraystretch}{1.3}
\setlength{\tabcolsep}{4pt}
\begin{tabular}{@{}p{2.0cm} p{10.5cm}@{}}
\toprule
Theme & Corresponding content in this chapter \\
\midrule
Translation & The propagation-operator group (\cref{prop:propagation_group}); the evolution semigroup (\cref{def:evolution_semigroup}); periodic closure of propagation (\cref{cor:propagation_period_closure}); additive translation of Bayesian updating in log-belief space (\cref{thm:bayesian_translation_structure}); finite-order translation of natural-parameter coordinates under a conjugate prior (\cref{thm:conjugate_exact_finite_order}) \\
Scaling & Wavelength scaling in a medium (\cref{prop:medium_wavelength_scaling}); the refractive-index scaling of the medium propagation operator (\cref{def:medium_propagation_operator}); the learning rate as the scaling factor of parametric evolution (\cref{def:learning_rate}) \\
Rotation & Consistency of the propagation operator and the phase operator (\cref{prop:phase_propagation}); the complex exponential plane-wave representation (\cref{def:vacuum_em_wave}); the phase rotation of the CEAF (\cref{prop:ceaf_phase_rotation}) \\
Conservation & Energy conservation of vacuum light propagation (\cref{thm:vacuum_energy_conservation}); flux conservation of a ray tube (\cref{prop:geometric_optics_interface}); monotone accumulation of the natural parameters under a conjugate prior (\cref{cor:natural_parameter_monotone}) \\
Symmetry & The propagation operator commutes with constant-coefficient differential operators (\cref{prop:propagation_commutes_translation}); invariance of Bloch waves under periodic translation (\cref{thm:bloch_floquet}); pointwise activation preserves translation invariance (\cref{cor:pointwise_preserves_translation}) \\
Closure & Periodic closure of propagation (\cref{cor:propagation_period_closure}); confinement of propagation in a photonic band gap (\cref{cor:photonic_bandgap}); convergence closure of parametric evolution (\cref{thm:parametric_evolution_closure}) \\
Period & Frequency-domain periodicity of the dispersion relation (\cref{prop:dispersion_period}); Bloch band structure (\cref{thm:bloch_floquet}); periodic structure of the cosine-annealing learning rate (\cref{prop:cosine_annealing_period}) \\
\bottomrule
\end{tabular}
\end{center}
\end{remark}

\section*{Seven-theme comparison table for this chapter}

\begin{center}
\renewcommand{\arraystretch}{1.25}
\begin{tabular}{@{}p{2.2cm} p{10.5cm}@{}}
\toprule
Theme & Corresponding content in this chapter \\
\midrule
Translation & The free-space propagation operator is the algebraic realisation of spatial translation (\cref{def:propagation_operator}), and the family of propagation operators forms a translation group (\cref{prop:propagation_group}); the time direction is carried by the evolution semigroup (\cref{def:evolution_semigroup}); Bayesian updating in log-belief space is a data-dependent additive translation (\cref{thm:bayesian_translation_structure}); parametric evolution is a translation in the gradient direction (\cref{def:parametric_evolution}) \\
Scaling & Wavelength scaling in a medium (\cref{prop:medium_wavelength_scaling}); the learning rate as the scaling factor of parametric evolution (\cref{def:learning_rate}, \cref{def:learning_rate_schedule}) \\
Rotation & Consistency of the propagation operator and the phase operator (\cref{prop:phase_propagation}); the complex exponential activation function behaves as a pure phase rotation under propagation (\cref{prop:ceaf_phase_rotation}); the Bloch--Floquet phase in a periodic medium (\cref{thm:bloch_floquet}) \\
Conservation & Energy conservation of vacuum light propagation (\cref{thm:vacuum_energy_conservation}); the isometry of the propagation operator, that is, conservation of the $L^2$ norm (\cref{prop:propagation_isometry}); causality (\cref{cor:causality_speed}) \\
Symmetry & The propagation operator commutes with constant-coefficient differential operators (\cref{prop:propagation_commutes_translation}); pointwise activation commutes with the propagation operator (\cref{prop:pointwise_activation_commutes}), and pointwise activation preserves translation invariance (\cref{cor:pointwise_preserves_translation}) \\
Closure & Periodic closure of propagation (\cref{cor:propagation_period_closure}); convergence closure of parametric evolution (\cref{thm:parametric_evolution_closure}) \\
Period & Bloch--Floquet propagation gives the periodic structure in a periodic medium (\cref{thm:bloch_floquet}) and the photonic band gap (\cref{cor:photonic_bandgap}); the relation between propagation convolution and periodic convolution (\cref{thm:propagation_periodic_convolution}); the exact finite-order evolution of Bayesian updating under a conjugate prior (\cref{thm:conjugate_exact_finite_order}); periodic structure of the cosine-annealing learning rate (\cref{prop:cosine_annealing_period}) \\
\bottomrule
\end{tabular}
\end{center}


\section*{Exercises}

\begin{enumerate}[leftmargin=3em, label=\textbf{Exercise \arabic*}]
    \item Using Definition \ref{def:optical_path_functional} and Definition \ref{def:fermat_principle},
    explain how Fermat's principle reduces a ray trajectory to a first-order stationary curve of the optical-path functional.
    Use Lemma \ref{lem:optical_path_variation} to write down the first variation of the optical-path functional,
    and explain how the boundary terms are handled in a variation with fixed endpoints.

    \item Prove the fundamental lemma of the calculus of variations (Lemma \ref{lem:fundamental_lemma_calculus_of_variations}):
    let $f\in C^0([s_0,s_1],\mathbb{R}^n)$; if for every $C^1$ vector function $\mathbf{h}$ satisfying
    $\mathbf{h}(s_0)=\mathbf{h}(s_1)=\mathbf{0}$ we have
    \[
    \int_{s_0}^{s_1}f(s)\cdot\mathbf{h}(s)\,ds=0,
    \]
    then $f\equiv\mathbf{0}$.
    Construct a proof by contradiction using a cut-off function.

    \item Prove the ray equation (Theorem \ref{thm:ray_equation}):
    \[
    \frac{d}{ds}\bigl(n(\gamma(s))\dot\gamma(s)\bigr)=\nabla n(\gamma(s)).
    \]
    Complete the proof in four steps---`expand the first variation, integrate by parts, set the boundary terms to zero,
    apply the fundamental lemma of the calculus of variations'---and
    explain the degeneration to straight-line propagation in a homogeneous medium (Corollary \ref{cor:straight_line_uniform}).

    \item Prove the eikonal equation (Corollary \ref{cor:eikonal_equation}):
    if there exists an optical path function $S$ such that $\dot\gamma=\nabla S/n$, then
    \[
    |\nabla S(\mathbf{r})|^2=n(\mathbf{r})^2.
    \]
    Use $|\dot\gamma|=1$ to complete the proof,
    and explain that the ray direction is perpendicular to the wavefront $S=\text{const}$.

    \item Prove Theorem \ref{thm:snell_consequence} (Snell's law):
    at a flat interface between two homogeneous media,
    \[
    n_1\sin\theta_1=n_2\sin\theta_2.
    \]
    Complete the proof by varying the point of incidence on the interface with Fermat's principle and
    setting the first variation to zero.
    Further explain that total internal reflection occurs when $n_1>n_2$ and the angle of incidence exceeds the critical angle.

    \item Prove Proposition \ref{prop:propagation_group} and Proposition \ref{prop:propagation_isometry}:
    \begin{enumerate}[label=(\roman*)]
        \item the family of propagation operators $\{P_{\mathbf{a}}\}$ forms a faithful representation of the translation group $\mathbb{R}^n$ under composition;
        \item the propagation operator $P_{\mathbf{a}}$ is an isometry on $L^2(\mathbb{R}^n)$,
              hence a unitary operator.
    \end{enumerate}
    Complete the proof using the change of variable $\mathbf{y}=\mathbf{x}-\mathbf{a}$.

    \item Prove Proposition \ref{prop:propagation_fourier}:
    in the Fourier domain the propagation operator is equivalent to multiplication by the phase factor,
    \[
    \widehat{P_{\mathbf{a}}f}(\mathbf{k})=e^{-i\mathbf{k}\cdot\mathbf{a}}\hat{f}(\mathbf{k}).
    \]
    Write out the direct computation of the Fourier transform,
    and prove Corollary \ref{cor:propagation_period_closure}:
    when the propagation distance equals one wavelength, $a=\lambda$, the plane wave returns to the phase it had at departure.

    \item Prove the energy conservation of vacuum light propagation (Theorem \ref{thm:vacuum_energy_conservation}):
    for solutions of the vacuum Maxwell equations,
    \[
    U(t)=\int_{\mathbb{R}^3}\Bigl(\frac{\varepsilon_0}{2}|\mathbf{E}|^2
    +\frac{1}{2\mu_0}|\mathbf{B}|^2\Bigr)d\mathbf{x}
    \]
    satisfies $U(t)=U(0)$.
    Explain how to write the Maxwell equations as $\partial_t\mathbf{u}=A\mathbf{u}$,
    and complete the proof using the antisymmetry of $A$ and Stone's theorem.

    \item Prove that the propagation kernel satisfies the evolution equation and the initial condition (Theorem \ref{thm:propagation_kernel_satisfies_pde}):
    the propagation kernel $G(\mathbf{x},t;\mathbf{x}_0)$ satisfies
    \[
    \partial_tG=\mathcal{L}G,\qquad
    G(\mathbf{x},0;\mathbf{x}_0)=\delta(\mathbf{x}-\mathbf{x}_0).
    \]
    Complete the proof using the definition of $G$ and the properties of the evolution operator $E_t=e^{t\mathcal{L}}$.

    \item Prove Lemma \ref{lem:gaussian_fourier_duality}:
    the Fourier transform of the Gaussian kernel
    \[
    K_\Delta(x)=\frac{1}{\sqrt{2\pi}\Delta}e^{-x^2/(2\Delta^2)}
    \]
    is again a Gaussian kernel:
    \[
    \hat{K}_\Delta(\omega)=e^{-\Delta^2\omega^2/2}.
    \]
    Complete the computation by completing the square in the Gaussian integral,
    and explain why the variances are interchanged, $\Delta\leftrightarrow1/\Delta$.

    \item Prove Theorem \ref{thm:propagation_periodic_convolution}:
    the periodic convolution of the single-period layer degenerates on the spatial translation group to propagation convolution,
    \[
    (f\circledast_{\text{prop}}G)(\mathbf{x})
    =\lim_{W\to\delta}(f\circledast g)(\mathbf{x}),
    \]
    where $g$ is taken to be the propagation kernel and $S^{(d)}$ the phase kernel $S^{(1)}$.
    Complete the proof starting from Definition \ref{def:periodic_convolution} and \ref{def:propagation_convolution}.

    \item Verify the Bloch--Floquet propagation theorem (Theorem \ref{thm:bloch_floquet}):
    let the linear evolution operator $\mathcal{L}$ have spatial period $\mathbf{T}$;
    then there exist Bloch wave solutions
    \[
    u_{\mathbf{k},n}(\mathbf{x},t)
    =e^{i(\mathbf{k}\cdot\mathbf{x}-\omega_n(\mathbf{k})t)}\phi_{\mathbf{k},n}(\mathbf{x}),
    \quad
    \phi_{\mathbf{k},n}(\mathbf{x}+\mathbf{T})=\phi_{\mathbf{k},n}(\mathbf{x}).
    \]
    Explain how the commutation of $\mathcal{L}$ with $P_{\mathbf{T}}$ leads to a simultaneous spectral decomposition,
    and explain the periodicity of the dispersion relation in the frequency domain (Proposition \ref{prop:dispersion_period}).

    \item Prove the translational structure of Bayesian updating (Theorem \ref{thm:bayesian_translation_structure}):
    in the log-belief space $\mathcal{F}(\Theta)$,
    \[
    \Lambda(p')=\Lambda(p)+\ell_D,
    \]
    where $\ell_D(\theta)=\log p(D\mid\theta)-\log Z_D$.
    Explain how Bayesian updating is equivalent to an additive translation in log-belief space,
    and point out the structural isomorphism between Bayesian updating and the propagation operator
    in Corollary \ref{cor:bayesian_propagation_isomorphism}.

    \item Prove the exact finite-order evolution of Bayesian updating under a conjugate prior (Theorem \ref{thm:conjugate_exact_finite_order}):
    let the likelihood be an exponential family and the prior a conjugate prior, and let the observation sequence be $x_1,\dots,x_n$;
    then the posterior natural parameters are
    \[
    \tau_n=\tau_0+\sum_{i=1}^nT(x_i),\qquad
    \nu_n=\nu_0+n.
    \]
    Complete the proof using Bayes' formula and the definition of the conjugate prior,
    and verify the Beta--binomial model (Example \ref{ex:beta_binomial_exact})
    and the Gaussian--Gaussian model (Example \ref{ex:gaussian_gaussian_exact}).

    \item Prove the commutativity of pointwise activation and the propagation operator (Proposition \ref{prop:pointwise_activation_commutes}):
    let $\varphi$ be a pointwise activation function and $P_{\mathbf{a}}$ a propagation operator; then
    \[
    P_{\mathbf{a}}\circ\varphi=\varphi\circ P_{\mathbf{a}}.
    \]
    Compute both sides for arbitrary $f$ and $\mathbf{x}$ and verify that they are equal,
    and prove Corollary \ref{cor:pointwise_preserves_translation}:
    pointwise activation preserves translation invariance.

    \item Prove Proposition \ref{prop:ceaf_phase_rotation}:
    under translation the CEAF behaves as a pure phase rotation,
    \[
    \mathcal{E}(x+a)=\sum_{k=1}^{N}A_ke^{i(\omega_kx+\phi_k+\omega_ka)}.
    \]
    Write out the factorisation of $e^{i(\omega_k(x+a)+\phi_k)}$,
    and prove the periodicity condition of Corollary \ref{cor:ceaf_period}.

    \item Prove Theorem \ref{thm:parametric_evolution_convergence}:
    let the quadratic loss be $\mathcal{L}(\theta)=\frac12(\theta-\theta^*)^{\mathsf T}H(\theta-\theta^*)$,
    with $H$ symmetric positive definite and eigenvalues $0<\lambda_1\le\dots\le\lambda_d$.
    Then for parametric evolution
    \[
    \theta_{t+1}=\theta_t-\eta\nabla\mathcal{L}(\theta_t)
    \]
    the error along the $i$-th eigendirection is
    \[
    e_t^{(i)}=(1-\eta\lambda_i)^te_0^{(i)},
    \]
    and the necessary and sufficient condition for convergence is $0<\eta<2/\lambda_d$.
    Write out the complete derivation by eigenbasis decomposition.

    \item Verify Example \ref{ex:parametric_evolution_numeric}:
    take $d=2$, $H=\operatorname{diag}(1,4)$, $\theta^*=(0,0)$ and $\theta_0=(1,1)$.
    \begin{enumerate}[label=(\roman*)]
        \item compute the convergence factors of the two eigendirections when $\eta=0.3$, and decide whether convergence occurs;
        \item compute the convergence factors of the two eigendirections when $\eta=0.6$, and decide whether convergence occurs.
    \end{enumerate}
    Explain why the convergence condition is determined by the largest eigenvalue $\lambda_d=4$.

    \item Prove Theorem \ref{thm:parametric_evolution_closure}:
    let $\mathcal{L}$ be an $L$-smooth, $\mu$-strongly convex function and take $\eta=1/L$; then
    \[
    \|\theta_t-\theta^*\|^2
    \le\Bigl(1-\frac{\mu}{L}\Bigr)^t\|\theta_0-\theta^*\|^2.
    \]
    Complete the proof using the descent lemma and strong convexity,
    and explain why parametric evolution converges exponentially---this is the convergence closure of parametric evolution.

    \item Synthesis: combining the ray equation (Theorem \ref{thm:ray_equation}),
    the propagation-operator group (Proposition \ref{prop:propagation_group}),
    Bloch--Floquet propagation (Theorem \ref{thm:bloch_floquet})
    and the translational structure of Bayesian updating (Theorem \ref{thm:bayesian_translation_structure}),
    explain how `propagation' and `evolution' are unified in this chapter as two realisations of one and the same translation-group action:
    \begin{enumerate}[label=(\roman*)]
        \item physical level: light propagation corresponds to the spatial translation group $\mathbb{R}^n$;
        \item temporal level: the evolution semigroup corresponds to the time translation group $\mathbb{R}_{\ge0}$;
        \item information level: Bayesian updating corresponds to a data-dependent translation;
        \item learning level: parametric evolution corresponds to a translation in the gradient direction.
    \end{enumerate}
    Point out for each realisation the translation amount, the generator and the group structure.
\end{enumerate}
